\chapter{Envelopes}
\index{envelope!BFJ}
\index{obstacle problem!non-Archimedean}
\index{orthogonality!dependency on envelopes}
\label{chap:bfj-envelopes}

An envelope projects an arbitrary obstacle onto the psh cone.  The
construction is one line long and its three basic theorems are not.  The
upper-semicontinuous regularization of a bounded family of psh functions is
psh; a \emph{continuous} obstacle has a \emph{continuous} envelope; and the
Monge--Amp\`ere measure of that envelope is concentrated on the contact set.
These three statements have three different engines.  The first runs on
skeletal equicontinuity and pure upper-envelope calculus.  The second runs on
asymptotic multiplier ideals, and it is the point at which residue
characteristic zero enters the original theory in an essential and
irremovable way.  The third is not proved in this chapter at all: it belongs
to the energy theory of \Cref{chap:bfj-differentiability}, and the present
chapter proves only its model-level shadow, in dimension one and in the
polyhedral convex category, where it is an elementary complementary-slackness
statement.

The reason for keeping the three apart is not fastidiousness.  Each is used
later under different hypotheses, and each fails, or is unknown, outside a
different region of the hypothesis matrix.  A reader who remembers only that
``envelopes behave well'' will, sooner or later, deduce the solution of the
non-Archimedean Monge--Amp\`ere equation over a field where no ingredient of
its proof is available.

\paragraph{Standing hypotheses.}
We retain the original BFJ hypotheses fixed at the beginning of
\Cref{chap:bfj-model-forms}: \(K\) is complete and discretely valued with
residue characteristic zero, \(v(\pi)=1\), \(X\) is smooth, connected and
projective of dimension \(n\) over \(K\), and \(\theta\) is a closed model
\((1,1)\)-form with ample generic class, as in
\eqref{eq:ch21-standing-theta}.  Every statement whose proof uses less than
this is marked; every later extension that needs a different hypothesis
signature is stated separately in \Cref{sec:ch23-later}, with its own date
and its own source.

\paragraph{What the reader should learn.}
By the end of the chapter the reader should be able to
manipulate \(P_\theta\) without silently inserting or deleting a
regularization star; recognize which of the two envelope definitions in
circulation is in force and why they agree here; reconstruct the base-ideal
proof of continuity, including the exact algebraic input and the exact place
where it is used; compute an envelope on a curve skeleton and on a polytope
by hand; and locate precisely where orthogonality is proved, what it costs,
and why no argument in this chapter delivers it.

\paragraph{Prerequisites.}
The compactness and equicontinuity architecture of
\Cref{chap:bfj-compactness}, in particular the skeletal estimate
\Cref{thm:bfj-skeletal-equicontinuity}; the definition and reconstruction
theorems of \Cref{chap:bfj-psh}; the upper-envelope, Dini, and Choquet
calculus of \Cref{sec:usc-envelopes} in \Cref{chap:topological-tools}; and,
for the algebraic input only, \Cref{sec:appC-asymptotic-multiplier} and
\Cref{sec:appC-global-generation} of \Cref{app:ideals}.

\section{Orientation}
\label{sec:ch23-orientation}

\subsection{Three operators that must not be confused}
\label{sec:ch23-three-operators}

Three distinct operations are all called ``taking an envelope'' in the
literature, and they are all denoted by some variant of \(P\).  We name them
once and never conflate them.

\begin{enumerate}
\item The \emph{obstacle envelope} \(P_\theta(f)\) of a function
\(f:X^{\An}\to\Rbar\): the pointwise supremum of all \(\theta\)-psh
functions lying below \(f\).  The input is an arbitrary function, the output
is a candidate psh function, and the defining inequality is a constraint
\emph{from above}.
\item The \emph{regularized supremum} \((\sup_\alpha\varphi_\alpha)^*\) of a
family of \(\theta\)-psh functions.  Here the input is already a family of
psh functions and the only issue is upper semicontinuity of the supremum.
\item The \emph{model obstacle envelope} \(P^{\mathcal D}_\theta(f)\),
obtained by restricting the competitors in (1) to \(\theta\)-psh
\emph{real model} functions in the convention of this chapter.
The rational-model envelope of
\cite[Definition~2.8]{GJKM19} is a distinct object, defined below using
numerical nefness and not a previously constructed singular psh class.
\end{enumerate}

Operations (1) and (2) are related but not identical: \(P_\theta(f)\) is the
supremum of a family that is closed under finite maxima and bounded above by
\(f\), so (2) applies to it, but the conclusion of (2) is about
\((P_\theta(f))^*\) and gives no information about \(P_\theta(f)\) itself
unless the obstacle is known to be upper semicontinuous.  Operations (1) and
(3) have different competitor sets; that
\(P^{\mathcal D}_\theta(f)\) and \(P_\theta(f)\) agree, in the present
hypothesis regime and for a continuous obstacle, is
\Cref{prop:ch23-model-vs-psh-envelope} below, and it is a theorem, not a
convention.

\begin{center}
\fbox{\begin{minipage}{0.91\textwidth}
\textbf{Notational rule for this chapter.}  The symbol \(P_\theta(f)\)
\emph{always} denotes the raw supremum of
\Cref{def:ch23-envelope}, never its regularization.  Whenever a
regularization is meant, the star is written.  The only statements in which
the star may be deleted are those in which upper semicontinuity of the raw
supremum has been proved; each such deletion is justified at the point where
it is made.
\end{minipage}}
\end{center}

\subsection{The hypothesis signature of each theorem}
\label{sec:ch23-signature}

The three principal results of the chapter carry the following signatures in
the sense of the hypothesis matrix of the conventions sheet.

\begin{center}
\small
\begin{tabular}{p{.26\textwidth}|p{.30\textwidth}|p{.34\textwidth}}
\textbf{Statement}&\textbf{Uses}&\textbf{Does not use}\\ \hline
\Cref{thm:ch23-upper-envelope}
 (supremum stability)
 &skeletal equicontinuity; usc calculus; Dini for nets
 &multiplier ideals; residue characteristic zero \emph{beyond} the
  availability of SNC models\\
\Cref{thm:bfj-envelope-regularity}
 (envelope continuity)
 &asymptotic multiplier ideals; subadditivity; uniform global generation;
  log resolutions
 &compactness; orthogonality; the Monge--Amp\`ere operator\\
\Cref{thm:ch23-model-orthogonality-curve}
 (model orthogonality, \(n=1\))
 &the linear description of relative nefness on a curve model
 &multiplier ideals; energy; differentiability of volumes\\
\end{tabular}
\end{center}

The middle row is the one that must be watched.  Residue characteristic zero
is used twice and in two different ways: once, weakly, to know that SNC
models are cofinal (\Cref{prop:snc-cofinal}), which is what makes the
skeletal definition of \Cref{def:ch21-bfj-psh} reasonable at all; and once,
essentially, to have asymptotic multiplier ideals with subadditivity and a
fixed-twist global generation theorem.  Removing the second use is the whole
content of the later literature discussed in \Cref{sec:ch23-later}, and it is
removed only at the cost of new hypotheses of a completely different kind.

\subsection{What is proved here and what is imported}
\label{sec:ch23-imported}

Everything in this chapter is proved in full from the following two imported
packages, and from the results of
\Cref{chap:bfj-model-forms,chap:bfj-psh,chap:bfj-compactness}.

\begin{enumerate}
\item \emph{Asymptotic multiplier ideals on a regular model:} their
existence, the containment \(\mathfrak a_m\subset\mathcal J(\mathfrak
a_\bullet^m)\), subadditivity
\(\mathcal J(\mathfrak a_\bullet^{m\ell})\subset\mathcal J(\mathfrak
a_\bullet^m)^\ell\), and a global generation statement with a twist
independent of \(m\).  These are collected as
\Cref{thm:ch23-uniform-generation-input}, with exact references, and their
construction is reviewed in \Cref{sec:appC-asymptotic-multiplier} and
\Cref{sec:appC-global-generation}.
\item \emph{Density of model functions:} \(\mathcal D(X)_{\mathbf Q}\) is
uniformly dense in \(C^0(X^{\An})\).  This is
\Cref{prop:metrics-model-function-density}, proved in \Cref{chap:metrics}
by a lattice Stone--Weierstrass argument.
\end{enumerate}

Nothing else is imported.  In particular, no result of
\Cref{chap:bfj-ma,chap:bfj-energy,chap:bfj-differentiability} is used
anywhere in this chapter, which is what makes the dependency chain
\eqref{eq:ch23-dependency-chain} non-circular.

\subsection{Plan}
\label{sec:ch23-plan}

\Cref{sec:ch23-obstacle-operator} sets up the obstacle operator, proves the
formal calculus, and compares the two definitions of the envelope.
\Cref{sec:ch23-upper-envelopes} proves supremum stability and the
usc-regularization statements, together with the fixed-skeleton countable
selection principle.  \Cref{sec:ch23-continuity} proves continuity of the
envelope of a continuous obstacle.  \Cref{sec:ch23-model-orthogonality}
proves orthogonality at model level in the two cases in which it is
elementary, and explains exactly why the argument stops there.
\Cref{sec:ch23-later} records the later continuity theorems with their exact
extra hypotheses.  \Cref{sec:ch23-worked} contains the computations, carried
to numbers, and \Cref{sec:ch23-dictionary} the chapter dictionary.

\section{The obstacle operator}
\label{sec:ch23-obstacle-operator}

\subsection{Definitions}
\label{sec:ch23-definitions}

\begin{definition}\label{def:ch23-envelope}
For an extended-real function \(f\) on \(X^{\An}\), define the raw envelope
\begin{equation}\label{eq:ch23-envelope}
 P_\theta(f)(x)
 :=\sup\{\varphi(x):
          \varphi\in\PSH(X,\theta),\ \varphi\leq f\text{ on }X^{\An}\}.
\end{equation}
If the admissible family is empty, the supremum is \(-\infty\).  For an
arbitrary obstacle, the \emph{psh envelope} means
\(P_\theta(f)^*\), the upper-semicontinuous regularization, unless the raw
supremum is already upper semicontinuous.
\end{definition}

Write
\begin{equation}\label{eq:ch23-admissible-family}
 \mathcal A_\theta(f)
 :=\{\varphi\in\PSH(X,\theta):\varphi\leq f\text{ on }X^{\An}\}
\end{equation}
for the admissible family, so that \(P_\theta(f)=\sup\mathcal A_\theta(f)\)
pointwise.  Three features of \eqref{eq:ch23-envelope} deserve to be stated
before anything is proved.

First, the constraint \(\varphi\leq f\) is imposed \emph{everywhere on
\(X^{\An}\)}, not merely on the skeleta.  For a psh competitor this is a
genuinely stronger requirement than its skeletal shadow, because a psh
function is determined by its skeletal restrictions
(\Cref{thm:bfj-skeletal-characterization}) while a general obstacle is not.
Second, \(\mathcal A_\theta(f)\) is stable under finite maxima, by
\Cref{prop:ch21-elementary-stability}, hence is a directed family; this is
what will eventually let Dini's theorem replace a bare supremum by a uniform
limit.  Third, no upper semicontinuity of \(f\) is assumed, and none is
inherited by \(P_\theta(f)\): the raw supremum of an upward-directed family
of upper semicontinuous functions need not be upper semicontinuous, and
\Cref{ex:ch23-bounded-obstacle-failure} exhibits a bounded obstacle for
which \(P_\theta(f)\) is not upper semicontinuous and for which its
regularization is not admissible.  This is exactly the phenomenon isolated in
\Cref{ex:sup-not-usc}.

\begin{definition}[Model obstacle envelope]
\label{def:ch23-model-envelope}
For \(f:X^{\An}\to\Rbar\), put
\begin{equation}\label{eq:ch23-model-envelope}
 P^{\mathcal D}_\theta(f)(x)
 :=\sup\{\varphi(x):
   \varphi\in\mathcal D(X)_{\mathbf R}\cap\PSH(X,\theta),\
   \varphi\leq f\text{ on }X^{\An}\},
\end{equation}
again with the convention that an empty supremum equals \(-\infty\).
\end{definition}

By \Cref{thm:ch21-model-consistency} a model function is \(\theta\)-psh in
the sense of \Cref{def:ch21-bfj-psh} if and only if \(\theta+\ddc\varphi\) is
a semipositive model form, so the competitor set in
\eqref{eq:ch23-model-envelope} is unambiguous and is contained in that of
\eqref{eq:ch23-envelope}.  Hence
\begin{equation}\label{eq:ch23-model-envelope-below}
 P^{\mathcal D}_\theta(f)\leq P_\theta(f)
 \qquad\text{pointwise on }X^{\An}.
\end{equation}
\Cref{prop:ch23-model-vs-psh-envelope} compares this real-model
envelope with the raw all-psh envelope in the present hypothesis regime.
For later arbitrary-field statements we use a separate notation.  Write
\(\mathcal D_{\mathbf Q}(X)\) for the rational model functions and set
\[
 P^{\mathcal D_{\mathbf Q}}_\theta(f)
 :=\sup\{\varphi\in\mathcal D_{\mathbf Q}(X):
       \theta+\ddc\varphi\text{ is a semipositive model form},\quad
       \varphi\leq f\},\qquad \sup\varnothing=-\infty.
\]
Here \(\theta\) may be a real closed model form.  This is the envelope
of \cite[\S\S2.2--2.8]{GJKM19}.  Its competitor set and that of
\eqref{eq:ch23-model-envelope} are not definitionally the same; neither
is identified with the raw all-psh envelope outside a stated comparison
theorem.

Finally we record the third operation of \Cref{sec:ch23-three-operators} as a
definition, since it is the one whose output requires a star.

\begin{definition}[Regularized supremum of a psh family]
\label{def:ch23-regularized-sup}
Let \((\varphi_\alpha)_{\alpha\in A}\) be a nonempty family in
\(\PSH(X,\theta)\) that is uniformly bounded above on \(X^{\An}\).  Its
\emph{regularized supremum} is
\begin{equation}\label{eq:ch23-regularized-sup}
 \Bigl(\sup_{\alpha\in A}\varphi_\alpha\Bigr)^{*},
\end{equation}
the least upper semicontinuous majorant \eqref{eq:usc-regularization} of the
pointwise supremum.
\end{definition}

\subsection{Nonemptiness and finiteness}
\label{sec:ch23-finiteness}

\begin{lemma}\label{lem:ch23-envelope-finiteness}
Let \(f:X^{\An}\to\Rbar\) be bounded below by a real constant on \(X^{\An}\).
Then \(\mathcal A_\theta(f)\) is nonempty and contains a model function, and
\begin{equation}\label{eq:ch23-envelope-two-sided}
 \inf_{X^{\An}}f-2\|\psi\|_\infty
 \ \leq\ P^{\mathcal D}_\theta(f)
 \ \leq\ P_\theta(f)\ \leq\ \sup_{X^{\An}}f ,
\end{equation}
where \(\psi\) is any model function with \(\theta+\ddc\psi\) semipositive.
In particular both envelopes are real valued and bounded whenever \(f\) is
bounded.  Conversely, if \(f\) is bounded above and
\(\mathcal A_\theta(f)\ne\varnothing\), then \(P_\theta(f)\) is real valued
at every quasimonomial point.
\end{lemma}

\begin{proof}
By \Cref{prop:ch20-positive-representative} there is a model function
\(\psi\), determined on a sufficiently high model, with
\(\theta':=\theta+\ddc\psi\) represented by a relatively ample class; in
particular \(\theta'\) is semipositive, so the constant \(0\) lies in
\(\PSH(X,\theta')\) and therefore \(\psi\in\PSH(X,\theta)\) by
\eqref{eq:ch21-reference-class-example}.  A model function is continuous on
the compact space \(X^{\An}\), hence bounded.  Put
\[
 \chi:=\psi+\Bigl(\inf_{X^{\An}}f-\sup_{X^{\An}}\psi\Bigr),
\]
which is \(\theta\)-psh because \(\PSH(X,\theta)\) is stable under adding
real constants (\Cref{prop:ch21-elementary-stability}).  Then
\(\chi\leq\inf f\leq f\), so \(\chi\) is an admissible model competitor, and
\[
 \chi\geq\inf_{X^{\An}}\psi+\inf_{X^{\An}}f-\sup_{X^{\An}}\psi
 \geq\inf_{X^{\An}}f-2\|\psi\|_\infty .
\]
This proves nonemptiness and the leftmost
inequality of \eqref{eq:ch23-envelope-two-sided}.  The middle inequality is
\eqref{eq:ch23-model-envelope-below}, and the rightmost holds because every
competitor is bounded above by \(f\).

For the converse, let \(\varphi\in\mathcal A_\theta(f)\) and let \(x\) be
quasimonomial, say \(x\in\Delta_{\mathcal X}\) for a determining SNC model.
Clause (2) of \Cref{def:ch21-bfj-psh} makes
\(\varphi|_{\Delta_{\mathcal X}}\) finite, so
\(P_\theta(f)(x)\geq\varphi(x)>-\infty\); and \(P_\theta(f)\leq\sup f<+\infty\).
\end{proof}

For a continuous obstacle the hypothesis of
\Cref{lem:ch23-envelope-finiteness} holds automatically, since \(X^{\An}\) is
compact.  We record the conclusion in the form in which it is used.

\begin{corollary}\label{cor:ch23-continuous-finite}
For \(f\in C^0(X^{\An})\), the functions \(P_\theta(f)\) and
\(P^{\mathcal D}_\theta(f)\) are real valued and uniformly bounded on
\(X^{\An}\), with
\(\inf f-2\|\psi\|_\infty\le P^{\mathcal D}_\theta(f)\le P_\theta(f)\le\sup f\).
\end{corollary}

\begin{proof}
Immediate from \Cref{lem:ch23-envelope-finiteness}, since a continuous
function on a compact space is bounded.
\end{proof}

The bound \eqref{eq:ch23-envelope-two-sided} is crude but it is uniform in
\(\theta\) on a compact set of classes, which is exactly what part (7) of
\Cref{prop:ch23-formal-envelope-calculus} needs; we return to this in
\Cref{lem:ch23-concave-lipschitz}.

\subsection{The formal calculus}
\label{sec:ch23-formal-calculus}

The next proposition is the algebra of the operator \(P_\theta\).  Each of
its clauses is formal, in the sense that it uses only
\Cref{def:ch23-envelope} and the elementary closure properties of
\(\PSH(X,\theta)\) established in \Cref{prop:ch21-elementary-stability}; none
uses compactness, multiplier ideals, or any continuity property of the
output.  We isolate first the one quantitative ingredient, a Lipschitz
estimate for concave functions, because it is used in clause (7) and is easy
to state incorrectly.

\begin{lemma}
\label{lem:ch23-concave-lipschitz}
Let \(V\) be a finite-dimensional normed real vector space, let
\(\tau_0\in V\), let \(r>0\), and let
\(h:B(\tau_0,2r)\to\mathbf R\) be concave with \(|h|\leq M\).  Then
\begin{equation}\label{eq:ch23-concave-lipschitz}
 |h(\tau)-h(\tau')|\leq\frac{2M}{r}\,\|\tau-\tau'\|,
 \qquad \tau,\tau'\in B(\tau_0,r).
\end{equation}
\end{lemma}

\begin{proof}
We may assume \(\tau\ne\tau'\).  Put \(d=\|\tau-\tau'\|\), so \(0<d<2r\), and
let \(e=(\tau'-\tau)/d\).  The point \(\tau'':=\tau'+re\) satisfies
\(\|\tau''-\tau_0\|\leq\|\tau'-\tau_0\|+r<2r\), so it lies in the larger
ball, and
\[
 \tau'=\frac{r}{d+r}\,\tau+\frac{d}{d+r}\,\tau''
\]
exhibits \(\tau'\) as a convex combination of \(\tau\) and \(\tau''\).
Concavity gives
\[
 h(\tau')\geq\frac{r}{d+r}h(\tau)+\frac{d}{d+r}h(\tau'').
\]
Multiplying by \((d+r)/r>0\) and rearranging,
\[
 h(\tau)-h(\tau')
 \leq\frac{d}{r}\bigl(h(\tau)-h(\tau'')\bigr)
 \leq\frac{2M}{r}\,d ,
\]
the last step by \(|h|\leq M\).  Exchanging \(\tau\) and \(\tau'\) gives
\eqref{eq:ch23-concave-lipschitz}.
\end{proof}

\begin{remark}
\label{rem:ch23-concave-lipschitz-sharp}
The passage to the concentric half-ball cannot be dropped: \(h(t)=\sqrt t\)
is concave and bounded on \((0,1)\) and is not Lipschitz there.  The constant
in \eqref{eq:ch23-concave-lipschitz} degenerates as \(r\downarrow0\), which
is why clause (7) below first fixes an open neighborhood of \(\theta\) on
which the two-sided bound of \Cref{lem:ch23-envelope-finiteness} is uniform,
and only afterwards lets \(\theta_j\) approach \(\theta\).
\end{remark}

\begin{proposition}\label{prop:ch23-formal-envelope-calculus}
Whenever the displayed envelopes are finite, the following identities and
inequalities hold.

\begin{enumerate}
\item If \(f\leq g\), then \(P_\theta(f)\leq P_\theta(g)\).
\item For \(c\in\mathbf R\),
\begin{equation}\label{eq:ch23-translation}
                         P_\theta(f+c)=P_\theta(f)+c.
\end{equation}

\item If \(0\leq t\leq1\), then
\begin{equation}\label{eq:ch23-joint-concavity}
 P_{t\theta+(1-t)\eta}(tf+(1-t)g)
 \geq tP_\theta(f)+(1-t)P_\eta(g).
\end{equation}

\item If \(u\) is a model function, then
\begin{equation}\label{eq:ch23-reference-change}
 P_{\theta+\ddc u}(f-u)=P_\theta(f)-u.
\end{equation}

\item If \(f,g\) are bounded, then
\begin{equation}\label{eq:ch23-envelope-contraction}
       \|P_\theta(f)-P_\theta(g)\|_\infty
       \leq\|f-g\|_\infty.
\end{equation}

\item If \(f\in\PSH(X,\theta)\), then \(P_\theta(f)=f\).
\item Let \(\mathcal X\) determine \(\theta\), and let
\(\theta_j\to\theta\) in \(N^1(\mathcal X/R)\), with ample generic classes.
For every fixed continuous \(f\),

\[
                       P_{\theta_j}(f)\longrightarrow P_\theta(f)
                       \quad\text{uniformly}.
\]
\end{enumerate}
\end{proposition}

\begin{proof}
\emph{(1) Monotonicity.}  If \(f\leq g\) then every \(\theta\)-psh
\(\varphi\leq f\) also satisfies \(\varphi\leq g\), so
\(\mathcal A_\theta(f)\subset\mathcal A_\theta(g)\) and the pointwise
suprema compare accordingly.

\emph{(2) Translation.}  The class \(\PSH(X,\theta)\) is stable under adding
a real constant (\Cref{prop:ch21-elementary-stability}), so
\(\varphi\mapsto\varphi+c\) is a bijection
\(\mathcal A_\theta(f)\to\mathcal A_\theta(f+c)\).  Taking suprema and using
that adding a finite constant commutes with a supremum in \(\Rbar\) gives
\eqref{eq:ch23-translation}.  The hypothesis \(c\in\mathbf R\) is not
decorative: for \(c=\pm\infty\) the right-hand side is not defined by the
conventions of \Cref{chap:topological-tools}.

\emph{(3) Joint concavity.}  Let \(\varphi\in\mathcal A_\theta(f)\) and
\(\psi\in\mathcal A_\eta(g)\), and let \(0\leq t\leq1\).  Convexity of
\(\PSH\) in the pair (potential, form) is the statement that
\(t\varphi+(1-t)\psi\in\PSH(X,t\theta+(1-t)\eta)\); it holds because upper
semicontinuity and the retraction inequality
\eqref{eq:ch21-retraction-inequality} are preserved by convex combinations,
because a convex combination of the corresponding uniform model approximants
on a fixed skeleton is again one, and because the nef cone is convex, which
is what makes those combinations \(t\theta+(1-t)\eta\)-psh at model level.
This is the content of \Cref{prop:ch21-elementary-stability} applied on a
common determining model.  Moreover
\(t\varphi+(1-t)\psi\leq tf+(1-t)g\).  Taking the supremum over
\(\varphi\) and then over \(\psi\), and using superadditivity of suprema,
gives \eqref{eq:ch23-joint-concavity}.  The inequality is genuinely one-sided:
already for \(\theta=\eta\) and \(g=f\) equality would assert that the
envelope is affine in the obstacle, which is false, as the two-obstacle
computation in \Cref{ex:ch23-curve-envelope-numbers} shows.

\emph{(4) Reference change.}  By
\eqref{eq:ch21-reference-class-example} the map \(\varphi\mapsto\varphi-u\)
is a bijection
\(\PSH(X,\theta)\to\PSH(X,\theta+\ddc u)\); it carries the constraint
\(\varphi\leq f\) to \(\varphi-u\leq f-u\), so it is a bijection
\(\mathcal A_\theta(f)\to\mathcal A_{\theta+\ddc u}(f-u)\).  Since \(u\) is
real valued, subtracting it commutes with the supremum, giving
\eqref{eq:ch23-reference-change}.

\emph{(5) Contraction.}  Put \(C=\|f-g\|_\infty<\infty\).  Then \(f\leq g+C\),
so (1) and (2) give \(P_\theta(f)\leq P_\theta(g)+C\).  Exchanging \(f\) and
\(g\) gives the opposite inequality, and the two together are
\eqref{eq:ch23-envelope-contraction}.  Note that the difference on the left
is formed only after \Cref{lem:ch23-envelope-finiteness} has made both terms
real valued; for unbounded obstacles the two-sided estimate must be written
without subtraction.

\emph{(6) Idempotence on the psh class.}  If \(f\in\PSH(X,\theta)\), then
\(f\in\mathcal A_\theta(f)\), so \(P_\theta(f)\geq f\); and every member of
\(\mathcal A_\theta(f)\) is \(\leq f\), so \(P_\theta(f)\leq f\).

\emph{(7) Continuity in the class.}  Fix a model \(\mathcal X\) determining
\(\theta\); after passing to a higher model,
\Cref{prop:ch20-positive-representative} provides a model function \(v\),
determined there, with \(\theta+\ddc v\) represented by a relatively ample
class \(\alpha\).  Ampleness is an open condition in the finite-dimensional
real vector space \(N^1(\mathcal X/R)\), so there is \(\rho>0\) such that
\[
 V:=\{\tau\in N^1(\mathcal X/R):\|\tau-\theta\|<4\rho\}
\]
consists of classes with \(\tau+\ddc v\) relatively ample.  For
\(\tau\in V\) the constant \(0\) is \((\tau+\ddc v)\)-psh, so \(v\) is
\(\tau\)-psh, and \Cref{lem:ch23-envelope-finiteness} gives the bounds
\begin{equation}\label{eq:ch23-uniform-class-bounds}
       \inf_{X^{\An}} f-2\|v\|_\infty
       \leq P_\tau(f)(x)\leq\sup_{X^{\An}} f,
       \qquad \tau\in V,\ x\in X^{\An},
\end{equation}
which are uniform in both \(\tau\in V\) and \(x\in X^{\An}\).  Let \(M\) be
the maximum of the absolute values of the two outer terms in
\eqref{eq:ch23-uniform-class-bounds}; it is finite because \(f\) is
continuous on a compact space, and it does not depend on \(x\) or \(\tau\).

Fix \(x\in X^{\An}\).  Taking \(f=g\) and \(\eta=\tau'\) in (3) shows that
\[
 h_x:V\longrightarrow\mathbf R,
 \qquad h_x(\tau)=P_\tau(f)(x)
\]
is concave, and \(|h_x|\leq M\) on \(V\), a ball of radius \(4\rho\)
about \(\theta\).  \Cref{lem:ch23-concave-lipschitz}, applied with
\(r=2\rho\), gives
\[
 |h_x(\tau)-h_x(\tau')|\leq\frac{M}{\rho}\|\tau-\tau'\|,
 \qquad\tau,\tau'\in B(\theta,2\rho),
\]
with a constant independent of \(x\).  Taking \(\tau'=\theta\), \(\tau=\theta_j\)
for \(j\) large enough that \(\|\theta_j-\theta\|<2\rho\), and then the
supremum over \(x\), yields
\[
       \|P_{\theta_j}(f)-P_\theta(f)\|_\infty
       \leq \frac{M}{\rho}\,\|\theta_j-\theta\|\longrightarrow0 ,
\]
which is (7).
\end{proof}

Two comments on the scope of the proposition are in order.  The contraction
estimate (5) is formal, but only after both raw envelopes are known to be
finite; that is \Cref{lem:ch23-envelope-finiteness}, and it is the sole place
where ampleness of \(\{\theta\}\) has been used so far.  Continuity of the
\emph{output} is a completely separate theorem and is proved only in
\Cref{sec:ch23-continuity}.  Nothing in
\Cref{prop:ch23-formal-envelope-calculus} asserts that \(P_\theta(f)\) is
\(\theta\)-psh, or even upper semicontinuous.

The list can be extended by the same methods; we record the additional
identities that are used later.

\begin{proposition}
\label{prop:ch23-further-calculus}
Let \(f,g:X^{\An}\to\Rbar\) be bounded, and assume
\(\mathcal A_\theta(f)\ne\varnothing\).  Then:
\begin{enumerate}
\item \emph{(Homogeneity)} for \(t>0\),
\begin{equation}\label{eq:ch23-homogeneity}
 P_{t\theta}(tf)=t\,P_\theta(f).
\end{equation}
\item \emph{(Idempotence)} \(P_\theta\bigl(P_\theta(f)\bigr)=P_\theta(f)\).
\item \emph{(Minima and maxima)}
\begin{equation}\label{eq:ch23-min-max}
 \begin{aligned}
 P_\theta(\min\{f,g\})&\leq\min\{P_\theta(f),P_\theta(g)\},\\
 P_\theta(\max\{f,g\})&\geq\max\{P_\theta(f),P_\theta(g)\}.
 \end{aligned}
\end{equation}
\item \emph{(Retraction majorant)} for every SNC model \(\mathcal X\) on
which \(\theta\) is determined,
\begin{equation}\label{eq:ch23-envelope-retraction-bound}
 P_\theta(f)\leq P_\theta(f)\circ p_{\mathcal X}
 \qquad\text{whenever }P_\theta(f)\in\PSH(X,\theta).
\end{equation}
\end{enumerate}
\end{proposition}

\begin{proof}
(1)  The map \(\varphi\mapsto t\varphi\) is a bijection from
\(\PSH(X,\theta)\) onto \(\PSH(X,t\theta)\) for \(t>0\): the retraction
inequality and upper semicontinuity are preserved, and a model function
\(\varphi_D\) is \(\theta\)-psh exactly when \(t\varphi_D=\varphi_{tD}\) is
\(t\theta\)-psh, since \(t(\theta+\ddc\varphi_D)=t\theta+\ddc(t\varphi_D)\)
and the nef cone is a cone.  Uniform skeletal approximants scale
accordingly.  The bijection carries \(\varphi\leq f\) to \(t\varphi\leq tf\),
and multiplying a supremum by \(t>0\) is legitimate in \(\Rbar\).

(2)  We have \(P_\theta(f)\leq f\) by definition of the admissible family
and the bound \(\varphi\le f\) for each competitor, so (1) of
\Cref{prop:ch23-formal-envelope-calculus} gives
\(P_\theta(P_\theta(f))\leq P_\theta(f)\).  Conversely every
\(\varphi\in\mathcal A_\theta(f)\) satisfies \(\varphi\leq P_\theta(f)\),
hence \(\varphi\in\mathcal A_\theta(P_\theta(f))\); taking suprema gives the
reverse inequality.

(3)  Both statements follow from monotonicity: \(\min\{f,g\}\leq f\) and
\(\min\{f,g\}\leq g\) give the first, and \(f\leq\max\{f,g\}\),
\(g\leq\max\{f,g\}\) the second.  Neither is an equality in general.  For
the first, let \(f\) and \(g\) be two \(\theta\)-psh model functions whose
restrictions to a one-dimensional face of some dual complex are affine with
different slopes.  Their minimum is concave and not affine near the
crossing, hence not convex there, hence not \(\theta\)-psh by clause (1) of
\Cref{thm:bfj-skeletal-characterization}.  By clause (6) of
\Cref{prop:ch23-formal-envelope-calculus} the right-hand side of the first
inequality in \eqref{eq:ch23-min-max} is \(\min\{f,g\}\) itself, while the
left-hand side is the envelope of a nonpsh obstacle and is therefore
strictly smaller somewhere.

(4)  This is \eqref{eq:ch21-retraction-inequality} applied to the psh
function \(P_\theta(f)\); the hypothesis is exactly what
\Cref{cor:ch23-continuous-obstacle-psh} supplies for continuous \(f\).
\end{proof}

\subsection{Retractions of a continuous obstacle}
\label{sec:ch23-retraction-uniform}

The following lemma is elementary but it is used at least four times below,
and it is the technical reason that continuous obstacles behave much better
than bounded ones.

\begin{lemma}
\label{lem:ch23-retraction-uniform}
Let \(f\in C^0(X^{\An})\).  Then
\begin{equation}\label{eq:ch23-retraction-uniform}
 \lim_{\mathcal X}\ \|f\circ p_{\mathcal X}-f\|_\infty=0,
\end{equation}
the limit being taken along the directed system of SNC models.  If
\(f\in\mathcal D(X)_{\mathbf R}\) is determined on \(\mathcal X_0\), then
\(f\circ p_{\mathcal X}=f\) for every \(\mathcal X\geq\mathcal X_0\).
\end{lemma}

\begin{proof}
The second assertion is the definition of a model function determined on
\(\mathcal X_0\): such a function is constant on the fibres of
\(p_{\mathcal X_0}\), and \(p_{\mathcal X_0}p_{\mathcal X}=p_{\mathcal X_0}\)
for \(\mathcal X\geq\mathcal X_0\) by
\eqref{eq:ch21-retraction-composition}, so
\(f\circ p_{\mathcal X}=f\circ p_{\mathcal X_0}\circ p_{\mathcal
X}=f\circ p_{\mathcal X_0}=f\).

For the first, let \(\varepsilon>0\).  By
\Cref{prop:metrics-model-function-density} there is a model function \(u\),
determined on some \(\mathcal X_0\), with \(\|f-u\|_\infty<\varepsilon/3\).
For every \(\mathcal X\geq\mathcal X_0\) and every \(x\in X^{\An}\),
\[
 |f(p_{\mathcal X}(x))-f(x)|
 \leq|f(p_{\mathcal X}(x))-u(p_{\mathcal X}(x))|
   +|u(p_{\mathcal X}(x))-u(x)|
   +|u(x)-f(x)| .
\]
The first and third terms are \(<\varepsilon/3\), and the middle term
vanishes by the second assertion.  Hence
\(\|f\circ p_{\mathcal X}-f\|_\infty<\varepsilon\) for all
\(\mathcal X\geq\mathcal X_0\), which is
\eqref{eq:ch23-retraction-uniform}.
\end{proof}

\begin{remark}
\label{rem:ch23-retraction-uniform-scope}
Uniformity in \eqref{eq:ch23-retraction-uniform} is what fails for a merely
bounded obstacle.  Pointwise one always has
\(f(p_{\mathcal X}(x))\to f(x)\) at every point of continuity of \(f\),
because \(p_{\mathcal X}(x)\to x\) in \(X^{\An}\) by
\Cref{thm:skeleton-inverse-limit}; but the convergence is a net convergence
in a space that is not metrizable, and no rate is available without density
of model functions.  Compare
\Cref{rem:net-vs-sequence-checklist}.
\end{remark}

\subsection{The two envelopes agree on quasimonomial points}
\label{sec:ch23-two-envelopes}

We can now settle the relation between \Cref{def:ch23-envelope} and
\Cref{def:ch23-model-envelope}.  The point of the proposition is that it is
proved \emph{before} envelope continuity, so that it may be used inside the
proof of \Cref{thm:ch23-base-ideal-approximation} without circularity.

\begin{proposition}\label{prop:ch23-model-vs-psh-envelope}
Let \(f\in C^0(X^{\An})\).  Then
\begin{equation}\label{eq:ch23-two-envelopes-quasimonomial}
 P^{\mathcal D}_\theta(f)(x)=P_\theta(f)(x)
 \qquad\text{for every quasimonomial }x\in X^{\An}.
\end{equation}
Consequently the two envelopes have the same restriction to every dual
complex, and
\begin{equation}\label{eq:ch23-two-envelopes-inf}
 \inf_{\mathcal X}\,P^{\mathcal D}_\theta(f)\circ p_{\mathcal X}
 =\inf_{\mathcal X}\,P_\theta(f)\circ p_{\mathcal X}
 \qquad\text{on }X^{\An}.
\end{equation}
\end{proposition}

\begin{proof}
By \eqref{eq:ch23-model-envelope-below} only the inequality
\(P_\theta(f)(x)\leq P^{\mathcal D}_\theta(f)(x)\) requires proof.  Fix a
quasimonomial \(x\), so \(x\in\Delta_{\mathcal Y}\) for some SNC model
\(\mathcal Y\) on which \(\theta\) is determined, and fix \(\varepsilon>0\).
By \Cref{cor:ch23-continuous-finite} the value \(P_\theta(f)(x)\) is finite,
so there is \(\psi\in\mathcal A_\theta(f)\) with
\begin{equation}\label{eq:ch23-competitor-near-x}
 \psi(x)>P_\theta(f)(x)-\varepsilon .
\end{equation}

\emph{Step 1: a model that sees the obstacle to within \(\varepsilon\).}
By \Cref{lem:ch23-retraction-uniform} choose an SNC model
\(\mathcal Z\geq\mathcal Y\) with
\begin{equation}\label{eq:ch23-obstacle-retraction-error}
 \|f\circ p_{\mathcal Z}-f\|_\infty<\varepsilon .
\end{equation}

\emph{Step 2: a model competitor on the skeleton of \(\mathcal Z\).}
By clause (2) of \Cref{def:ch21-bfj-psh} the restriction
\(\psi|_{\Delta_{\mathcal Z}}\) is a uniform limit of restrictions of
\(\theta\)-psh model functions.  Choose a \(\theta\)-psh model function
\(\chi\) with
\begin{equation}\label{eq:ch23-skeletal-approximation-error}
 \bigl|\chi-\psi\bigr|<\varepsilon
 \qquad\text{on }\Delta_{\mathcal Z}.
\end{equation}
Enlarging \(\mathcal Z\) if necessary, we may and do assume that \(\chi\) is
determined on \(\mathcal Z\); this does not disturb
\eqref{eq:ch23-obstacle-retraction-error}, which persists on higher models,
and it does not disturb \eqref{eq:ch23-skeletal-approximation-error} if the
approximation of Step~2 is performed after the enlargement.

\emph{Step 3: pushing the inequality off the skeleton.}  Since \(\chi\) is
determined on \(\mathcal Z\) we have \(\chi=\chi\circ p_{\mathcal Z}\), by
\Cref{lem:ch21-model-skeletal-properties}.  For every \(x'\in X^{\An}\),
\[
 \chi(x')=\chi(p_{\mathcal Z}(x'))
 \overset{\eqref{eq:ch23-skeletal-approximation-error}}{<}
 \psi(p_{\mathcal Z}(x'))+\varepsilon
 \leq f(p_{\mathcal Z}(x'))+\varepsilon
 \overset{\eqref{eq:ch23-obstacle-retraction-error}}{<}
 f(x')+2\varepsilon ,
\]
where the middle inequality is \(\psi\leq f\) evaluated at
\(p_{\mathcal Z}(x')\).  Hence
\(\chi-2\varepsilon\leq f\) on all of \(X^{\An}\), and
\(\chi-2\varepsilon\) is a \(\theta\)-psh model function, so it belongs to
the competitor set of \eqref{eq:ch23-model-envelope}.

\emph{Step 4: conclusion.}  Since \(x\in\Delta_{\mathcal Y}\subset
\Delta_{\mathcal Z}\), \eqref{eq:ch23-skeletal-approximation-error} and
\eqref{eq:ch23-competitor-near-x} give
\[
 P^{\mathcal D}_\theta(f)(x)
 \geq\chi(x)-2\varepsilon
 >\psi(x)-3\varepsilon
 >P_\theta(f)(x)-4\varepsilon .
\]
Letting \(\varepsilon\downarrow0\) proves
\eqref{eq:ch23-two-envelopes-quasimonomial}.

The remaining assertions are immediate: every point of a dual complex is
quasimonomial, and \(p_{\mathcal X}(x)\) is quasimonomial for every \(x\) and
every \(\mathcal X\), so the two functions in
\eqref{eq:ch23-two-envelopes-inf} are infima of the same family of
functions.
\end{proof}

\begin{remark}[Why \eqref{eq:ch23-two-envelopes-quasimonomial} is not an
equality everywhere, a priori]
\label{rem:ch23-two-envelopes-scope}
Nothing in the proof compares the two envelopes at a point that is not
quasimonomial.  After \Cref{thm:bfj-envelope-regularity} both functions are
continuous and agree on the dense set of quasimonomial points
(\Cref{thm:skeleton-inverse-limit}), hence agree everywhere; but that is a
consequence of continuity, not a substitute for it.  The corresponding
statement in \cite[Definition~2.8 and the discussion after it]{GJKM19} is
made in exactly this order, and in the positive-characteristic setting of
\Cref{sec:ch23-later} the identification of the two definitions is again
deduced from the continuity theorem, never assumed.
\end{remark}

\section{Upper envelopes of psh functions}
\label{sec:ch23-upper-envelopes}

The first analytic theorem concerns a family of psh functions, not a family
of model functions.  Its proof uses no algebraic geometry beyond the
equicontinuity estimate of \Cref{thm:bfj-skeletal-equicontinuity}: everything
else is the upper-envelope and Dini calculus of
\Cref{chap:topological-tools}.

\begin{theorem}\label{thm:ch23-upper-envelope}
Let \((\varphi_\alpha)_{\alpha\in A}\) be a nonempty family in
\(\PSH(X,\theta)\), uniformly bounded above on \(X^{\An}\), and set
\begin{equation}\label{eq:ch23-raw-psh-sup}
                         \varphi=\sup_\alpha\varphi_\alpha.
\end{equation}
Then \(\varphi^*\in\PSH(X,\theta)\),
\(\varphi^*=\varphi\) at every quasimonomial point, and
\begin{equation}\label{eq:ch23-upper-envelope-reconstruction}
 \varphi^*
   =\inf_{\mathcal X}\varphi\circ p_{\mathcal X}
   =\lim_{\mathcal X}\varphi\circ p_{\mathcal X}.
\end{equation}
\end{theorem}

\begin{proof}
The proof is in five steps.  Steps~1--3 take place on one fixed dual
complex and are pure analysis; Steps~4--5 assemble the skeletal information
into a global statement.

\emph{Step 1: reduction to an increasing net.}  Let \(\mathcal P(A)_{\mathrm
f}\) be the set of nonempty finite subsets of \(A\), directed by inclusion
as in \Cref{ex:directed-sets-of-the-book}, and put
\begin{equation}\label{eq:ch23-finite-maxima}
                  \varphi_F:=\max_{\alpha\in F}\varphi_\alpha,
                  \qquad F\in\mathcal P(A)_{\mathrm f}.
\end{equation}
Each \(\varphi_F\) lies in \(\PSH(X,\theta)\) by stability under finite
maxima (\Cref{prop:ch21-elementary-stability}).  The net
\((\varphi_F)_F\) is increasing, and
\(\sup_F\varphi_F=\sup_\alpha\varphi_\alpha=\varphi\) pointwise, since every
\(\alpha\) occurs in some \(F\) and every \(\varphi_F\) is a maximum of
members of the family.  Note that \(\mathcal P(A)_{\mathrm f}\) admits no
cofinal sequence when \(A\) is uncountable
(\Cref{ex:no-cofinal-sequence}); everything below is therefore formulated
for nets.

\emph{Step 2: equicontinuity on a fixed skeleton.}  Fix an SNC model
\(\mathcal X\) on which \(\theta\) is determined, write
\(\Delta=\Delta_{\mathcal X}\), and let \(d\) be the associated path metric:
\(d(x,y)\) is the infimum of the lengths of paths in \(\Delta\) from \(x\) to
\(y\), lengths being measured face by face in the fixed Euclidean structure.
This infimum is finite because \(\Delta\) is connected, the special fibre
being connected.  Let \(C_0\) be a uniform bound with
\(\varphi_\alpha\leq C_0\) on \(X^{\An}\) for all \(\alpha\).  Clause (1) of
\Cref{thm:bfj-skeletal-equicontinuity} gives a constant
\(L=L(\mathcal X,\theta)\), independent of the function, such that the
restriction to \(\Delta\) of \emph{every} \(\theta\)-psh function is
\(L\)-Lipschitz on each face; subdividing a path into segments each
contained in one face shows that it is then \(L\)-Lipschitz for \(d\).  Set
\(\omega(t)=Lt\).  A maximum of finitely many functions with modulus
\(\omega\) again has modulus \(\omega\), so
\begin{equation}\label{eq:ch23-common-modulus}
 \varphi_F(x)\leq\varphi_F(y)+\omega(d(x,y)),
 \qquad x,y\in\Delta,
\end{equation}
for every \(F\).

\emph{Step 3: the supremum is continuous on \(\Delta\), and the convergence
is uniform.}  The pointwise limit \(\varphi|_\Delta\) is finite: it is
bounded above by \(C_0\), and bounded below by
\(\varphi_{\alpha_0}|_\Delta\) for any fixed \(\alpha_0\), which is finite
by clause (1) of \Cref{thm:bfj-skeletal-characterization}.  Taking the
supremum over \(F\) in \eqref{eq:ch23-common-modulus} gives
\(\varphi(x)\leq\varphi(y)+\omega(d(x,y))\); interchanging \(x\) and \(y\)
gives the same bound with the roles reversed, so \(\varphi|_\Delta\) has the
modulus \(\omega\) and in particular is continuous.  The net of continuous
functions \(\varphi_F|_\Delta\) increases to the continuous function
\(\varphi|_\Delta\) on the compact space \(\Delta\), so
\Cref{prop:dini-nets}, applied to the decreasing net
\(\varphi|_\Delta-\varphi_F|_\Delta\) of continuous nonnegative functions,
gives
\begin{equation}\label{eq:ch23-dini-on-skeleton}
 \lim_F\ \bigl\|\varphi_F|_\Delta-\varphi|_\Delta\bigr\|_\infty=0 .
\end{equation}
Consequently \(\varphi\circ p_{\mathcal X}\) is continuous on \(X^{\An}\),
being the composition of a continuous function on \(\Delta\) with the
continuous retraction \(p_{\mathcal X}\)
(\Cref{prop:skeleton-retraction}).

\emph{Step 4: the reconstruction formula.}  If
\(\mathcal X'\geq\mathcal X\), then
\eqref{eq:ch21-retraction-composition} gives
\(p_{\mathcal X}p_{\mathcal X'}=p_{\mathcal X}\), and applying
\eqref{eq:ch21-retraction-inequality} to each \(\varphi_\alpha\) at the point
\(p_{\mathcal X'}(x)\) and taking suprema yields
\[
 \varphi(x)\leq\varphi(p_{\mathcal X'}(x))\leq\varphi(p_{\mathcal X}(x)),
\]
exactly as in \eqref{eq:ch21-decreasing-retractions}.  Thus
\((\varphi\circ p_{\mathcal X})_{\mathcal X}\) is a decreasing net of
continuous functions; put
\[
                        h:=\inf_{\mathcal X}
                             \varphi\circ p_{\mathcal X}
                        =\lim_{\mathcal X}\varphi\circ p_{\mathcal X},
\]
the second equality by \Cref{lem:monotone-net-converges}.  As an infimum of
continuous functions, \(h\) is upper semicontinuous
(\Cref{lem:usc-inf}); and \(h\geq\varphi\) by the display above.  Since
\(\varphi^*\) is the \emph{least} upper semicontinuous majorant of
\(\varphi\) (\Cref{prop:usc-majorant}), we get \(h\geq\varphi^*\).

For the reverse inequality, fix \(x\in X^{\An}\).  The net
\((p_{\mathcal X}(x))_{\mathcal X}\) converges to \(x\) by
\Cref{thm:skeleton-inverse-limit}, so upper semicontinuity of
\(\varphi^*\), in the net form \Cref{prop:net-semicontinuity}, gives
\begin{equation}\label{eq:ch23-usc-reconstruction-bound}
 \varphi^*(x)
 \geq\varlimsup_{\mathcal X}
           \varphi^*(p_{\mathcal X}(x))
 \geq\varlimsup_{\mathcal X}
           \varphi(p_{\mathcal X}(x))
 =h(x),
\end{equation}
the middle step because \(\varphi\leq\varphi^*\) and the last because the
net is decreasing, so its upper limit is its infimum.  This proves
\eqref{eq:ch23-upper-envelope-reconstruction}.

\emph{Step 5: \(\varphi^*\) is \(\theta\)-psh.}  It is upper semicontinuous
by construction and is not identically \(-\infty\), since
\(\varphi^*\geq\varphi\geq\varphi_{\alpha_0}\).  If
\(\mathcal Y\leq\mathcal X\), then
\(p_{\mathcal X}p_{\mathcal Y}=p_{\mathcal Y}\) and hence
\(\varphi\circ p_{\mathcal X}\circ p_{\mathcal Y}
=\varphi\circ p_{\mathcal Y}\) for every \(\mathcal X\geq\mathcal Y\);
taking the infimum over \(\mathcal X\geq\mathcal Y\), which is cofinal, in
\eqref{eq:ch23-upper-envelope-reconstruction} gives
\begin{equation}\label{eq:ch23-equality-on-skeleton}
                   \varphi^*\circ p_{\mathcal Y}
                   =\varphi\circ p_{\mathcal Y}.
\end{equation}
Evaluating \eqref{eq:ch23-equality-on-skeleton} at a point of
\(\Delta_{\mathcal Y}\), where \(p_{\mathcal Y}\) is the identity, shows
\(\varphi^*=\varphi\) on \(\Delta_{\mathcal Y}\), hence at every
quasimonomial point, since every such point lies on some dual complex.

The retraction inequality for \(\varphi^*\) now reads
\(\varphi^*\leq\varphi\circ p_{\mathcal Y}=\varphi^*\circ p_{\mathcal Y}\),
which is \eqref{eq:ch23-upper-envelope-reconstruction} combined with
\eqref{eq:ch23-equality-on-skeleton}; this is clause (1) of
\Cref{def:ch21-bfj-psh}.  For clause (2), fix \(\mathcal Y\) and use
\eqref{eq:ch23-equality-on-skeleton} together with
\eqref{eq:ch23-dini-on-skeleton}: the restriction
\(\varphi^*|_{\Delta_{\mathcal Y}}=\varphi|_{\Delta_{\mathcal Y}}\) is a
uniform limit of the restrictions \(\varphi_F|_{\Delta_{\mathcal Y}}\), each
of which is itself a uniform limit of restrictions of \(\theta\)-psh model
functions because \(\varphi_F\in\PSH(X,\theta)\).  A uniform limit of
elements of a uniformly closed set lies in that set, so
\(\varphi^*|_{\Delta_{\mathcal Y}}\) belongs to the uniform closure of the
restrictions of \(\theta\)-psh model functions.  Both clauses of
\Cref{def:ch21-bfj-psh} hold, and \(\varphi^*\in\PSH(X,\theta)\).
\end{proof}

\begin{remark}[Where each hypothesis is used]
\label{rem:ch23-upper-envelope-hypotheses}
Nonemptiness of the family is used twice: to know that \(\varphi\) is not
identically \(-\infty\), and to get a finite lower bound on each skeleton in
Step~3.  The uniform upper bound is used only in Step~3, but it cannot be
weakened to a pointwise bound: without it the increasing limit on a skeleton
may be \(+\infty\) at some point, and then no modulus of continuity is
available.  Stability under finite maxima
(\Cref{prop:ch21-elementary-stability}) is what makes the reduction of
Step~1 legitimate; it is a theorem about the nef cone
(\Cref{prop:ch20-model-max}), not a formality.
\end{remark}

\begin{corollary}
\label{cor:ch23-increasing-net}
Let \((\varphi_i)_{i\in I}\) be an increasing net in \(\PSH(X,\theta)\),
uniformly bounded above, and let \(\varphi=\sup_i\varphi_i\).  Then
\(\varphi^*\in\PSH(X,\theta)\), \(\varphi^*=\varphi\) on the union of all
dual complexes, and
\(\varphi^*=\varphi\) on \(X^{\An}\) if and only if \(\varphi\) is upper
semicontinuous.  If moreover \(\varphi\) is continuous, then
\(\varphi_i\to\varphi\) uniformly on every dual complex.
\end{corollary}

\begin{proof}
The first two assertions are \Cref{thm:ch23-upper-envelope} together with
the last assertion of \Cref{prop:usc-majorant}, which says that
\(\varphi=\varphi^*\) exactly when \(\varphi\) is upper semicontinuous.  The
last assertion is \eqref{eq:ch23-dini-on-skeleton}, since an increasing net
is cofinal in the net of its own finite maxima.
\end{proof}

The regularization in \Cref{thm:ch23-upper-envelope} cannot be omitted, and
the failure is not confined to pathological families.

\begin{example}[A bounded increasing family whose supremum is not upper
semicontinuous]
\label{ex:ch23-sup-not-usc}
Let \(X=\mathbf P^1_K\) with homogeneous coordinate \(z\), let
\(\mathcal X=\mathbf P^1_R\), and let \(\theta=c_1(\cO(1),\|\cdot\|_0)\) for
the model metric of \(\cO_{\mathbf P^1_R}(1)\).  Put
\begin{equation}\label{eq:ch23-log-pole-potential}
 \psi(x):=-\max\{v_x(z),0\}
 =\min\{\log|z(x)|,0\},
 \qquad x\in(\mathbf P^1_K)^{\An}.
\end{equation}
Then \(\psi\) is \(\theta\)-psh: it is upper semicontinuous because
\(x\mapsto|z(x)|\) is continuous; it satisfies the retraction inequality
because \(|z|\) does, by \Cref{lem:retraction-minorant}; and on the skeleton
of any SNC model its restriction is the model function of an anti-effective
vertical divisor whose curvature is nef, as verified in
\Cref{ex:ch23-pole-model-check}.  Its polar set is the single rigid point
\(z=0\), and \(\psi\leq0\) with \(\psi=0\) on \(\{|z|\geq1\}\).

For \(0<\varepsilon\leq1\) put \(\varphi_\varepsilon:=\varepsilon\psi\).
Each \(\varphi_\varepsilon\) is \(\theta\)-psh, being the convex combination
\(\varepsilon\psi+(1-\varepsilon)\cdot0\) of two \(\theta\)-psh functions,
and the family is increasing as \(\varepsilon\downarrow0\) and bounded above
by \(0\).  Its pointwise supremum is
\[
 \varphi=\sup_{0<\varepsilon\leq1}\varepsilon\psi
 =\begin{cases}
   0,&\psi(x)>-\infty,\\
   -\infty,&x=[0:1],
 \end{cases}
\]
because \(\varepsilon\psi(x)\to0\) for finite \(\psi(x)\) while
\(\varepsilon\cdot(-\infty)=-\infty\) for every \(\varepsilon>0\).  Thus
\(\varphi\) is not upper semicontinuous at \(z=0\): its value there is
\(-\infty\), while it equals \(0\) on the punctured neighborhood.  Its
regularization is \(\varphi^*\equiv0\), which is indeed \(\theta\)-psh, in
accordance with \Cref{thm:ch23-upper-envelope}; and \(\varphi^*=\varphi\) at
every quasimonomial point, since no quasimonomial point is a rigid point.
\end{example}

\begin{example}[The potential \eqref{eq:ch23-log-pole-potential} is
\(\theta\)-psh, with the model check written out]
\label{ex:ch23-pole-model-check}
Write \([Z:W]\) for the homogeneous coordinates, so that \(z=Z/W\), and let
\(\|\cdot\|_0\) be the model metric of \(\cO_{\mathbf P^1_R}(1)\).  In the
chart \(W\ne0\) one has \(\|Z\|_0=|z|/\max\{|z|,1\}\), whence
\begin{equation}\label{eq:ch23-section-weight}
 \log\|Z\|_0=\min\{\log|z|,0\}=\psi .
\end{equation}
For \(m\geq1\) let \(\mathfrak b_m\subset\cO_{\mathbf P^1_R}\) be the image
of
\[
 \cO(-1)\oplus\cO
 \xrightarrow{\ (Z,\ \pi^m)\ }
 \cO_{\mathbf P^1_R}.
\]
Its zero locus is the reduction of the point \(z=0\), which lies in the
special fibre, so \(\mathfrak b_m\) is a vertical ideal, and by
\eqref{eq:ch23-section-weight}
\begin{equation}\label{eq:ch23-truncated-pole}
 \log|\mathfrak b_m|=\max\{\psi,-m\}=:\psi_m .
\end{equation}
The sheaf \(\cO(1)\otimes\mathfrak b_m\) is generated by the global sections
\(Z\) and \(\pi^mW\), hence globally generated, so
\Cref{lem:ch20-generated-ideal-psh} makes \(\psi_m\) a \(\theta\)-psh model
function.  The net \((\psi_m)_{m\geq1}\) is decreasing with pointwise
infimum \(\psi\), and \(\psi\not\equiv-\infty\); by
\Cref{thm:ch21-decreasing-stability}, \(\psi\in\PSH(X,\theta)\).  Note also
that on any fixed dual complex \(\Delta\) the function \(\psi\) is finite and
continuous, so \(\psi_m=\psi\) on \(\Delta\) as soon as
\(m\geq-\min_\Delta\psi\): the truncation is invisible on each skeleton, and
only the rigid point \(z=0\) ever sees it.

It is worth following where the curvature mass goes.  Let
\(\mathcal X_m\to\mathbf P^1_R\) be the model obtained by \(m\) successive
blowups, the first at the reduction of \(z=0\) and each subsequent one at the
point of the last exceptional divisor lying over it.  Its special fibre is a
chain \(E_0\cup\cdots\cup E_m\) of rational curves with all multiplicities
equal to \(1\), with \((E_i\cdot E_{i+1})=1\), all other products zero, and
\[
 (E_0^2)=(E_m^2)=-1,\qquad (E_i^2)=-2\quad(1\leq i\leq m-1).
\]
The divisorial points satisfy \(v_{e_i}(z)=i\), so
\(\psi(e_i)=-i\) and \(\psi|_{\Delta_{\mathcal X_m}}=\varphi_{D_m}\) with
\(D_m=-\sum_{i=1}^m iE_i\) anti-effective, in agreement with
\eqref{eq:ch20-sign}.  With \(d_i=(\pi_m^*\cO(1)\cdot E_i)\), so that
\(d_0=1\) and \(d_i=0\) for \(i\geq1\), the test
\eqref{eq:ch20-curve-nef-test} gives the residual degrees
\[
 \begin{aligned}
 A_0&=1+(-1-0)=0,\\
 A_i&=(-(i-1)+i)+(-(i+1)+i)=0\quad(1\leq i\leq m-1),\\
 A_m&=(-(m-1)+m)=1.
 \end{aligned}
\]
All are nonnegative, confirming nefness directly, and the model measure of
\(\theta+\ddc\psi_m\) is the single atom \(\delta_{e_m}\), of total mass
\(1=\deg\cO(1)\).  As \(m\to\infty\) the atom migrates along the chain
towards the rigid point \(z=0\); in the limit the curvature of the singular
metric \(\|\cdot\|_0\rme^{-\psi}\) is the Dirac mass at that point, which is
the Poincar\'e--Lelong statement for the section \(Z\).
\end{example}

\begin{corollary}\label{cor:ch23-continuous-obstacle-psh}
If \(f\in C^0(X^{\An})\), then the raw envelope \(P_\theta(f)\) is already
upper semicontinuous and \(\theta\)-psh.  It is the largest
\(\theta\)-psh function dominated by \(f\).
\end{corollary}

\begin{proof}
By \Cref{cor:ch23-continuous-finite} the admissible family
\(\mathcal A_\theta(f)\) is nonempty and uniformly bounded above by
\(\sup f<\infty\).  \Cref{thm:ch23-upper-envelope} therefore applies and
shows \(P_\theta(f)^*\in\PSH(X,\theta)\).

We claim \(P_\theta(f)^*\leq f\).  Indeed \(P_\theta(f)\leq f\) by
definition, \(f\) is upper semicontinuous because it is continuous, and
\(P_\theta(f)^*\) is the least upper semicontinuous majorant of
\(P_\theta(f)\) by \Cref{prop:usc-majorant}; hence
\(P_\theta(f)^*\leq f\).  This is exactly the criterion in clause (6) of
\Cref{lem:usc-star-calculus}, applied to the family
\(\mathcal A_\theta(f)\) and the usc majorant \(f\).

Consequently \(P_\theta(f)^*\in\mathcal A_\theta(f)\), so
\(P_\theta(f)^*\leq\sup\mathcal A_\theta(f)=P_\theta(f)\).  Combined with
\(P_\theta(f)\leq P_\theta(f)^*\) this gives
\[
               P_\theta(f)^*\leq P_\theta(f)\leq P_\theta(f)^*,
\]
so the raw envelope equals its own regularization and is therefore upper
semicontinuous and \(\theta\)-psh.  Maximality is the statement that
\(P_\theta(f)\) is itself a member of \(\mathcal A_\theta(f)\) dominating
every other member, which we have just proved.
\end{proof}

The only property of \(f\) used in the proof is upper semicontinuity, so we
record the sharper statement.

\begin{corollary}
\label{cor:ch23-usc-obstacle}
Let \(f:X^{\An}\to\mathbf R\cup\{-\infty\}\) be upper semicontinuous and
bounded above, with \(\mathcal A_\theta(f)\ne\varnothing\).  Then
\(P_\theta(f)\in\PSH(X,\theta)\) and it is the largest \(\theta\)-psh
function dominated by \(f\).
\end{corollary}

\begin{proof}
Repeat the previous proof verbatim, replacing ``continuous'' by ``upper
semicontinuous'' and using \Cref{lem:ch23-envelope-finiteness} in place of
\Cref{cor:ch23-continuous-finite}.
\end{proof}

Semicontinuity of the obstacle from \emph{above} is what makes the argument
work, and it cannot be replaced by boundedness.

\begin{example}[A bounded obstacle whose envelope is not upper
semicontinuous]
\label{ex:ch23-bounded-obstacle-failure}
Keep the notation of \Cref{ex:ch23-sup-not-usc}: \(X=\mathbf P^1_K\),
\(\theta=c_1(\cO(1),\|\cdot\|_0)\), \(x_0=[0:1]\) the rigid point \(z=0\),
and \(\psi\) as in \eqref{eq:ch23-log-pole-potential}.  Define a bounded
obstacle by
\[
 f:=-\one_{\{x_0\}},
 \qquad\text{that is, } f(x_0)=-1 \text{ and } f=0 \text{ elsewhere}.
\]
Then \(f\) is bounded but not upper semicontinuous at \(x_0\).

Since \(\theta\) is represented by the relatively ample class of
\(\cO_{\mathbf P^1_R}(1)\), constants are \(\theta\)-psh; and
\(\PSH(X,\theta)\) is stable under finite maxima by
\Cref{prop:ch21-elementary-stability}.  For \(\varepsilon\in(0,1]\) put
\[
 \chi_\varepsilon:=\max\{\varepsilon\psi,\,-1\}\in\PSH(X,\theta).
\]
Then \(\chi_\varepsilon\leq0\) everywhere and
\(\chi_\varepsilon(x_0)=\max\{-\infty,-1\}=-1=f(x_0)\), so
\(\chi_\varepsilon\in\mathcal A_\theta(f)\).  For \(x\ne x_0\) the value
\(\psi(x)\) is finite, so
\(\chi_\varepsilon(x)\to\max\{0,-1\}=0\) as \(\varepsilon\downarrow0\).
Since every competitor satisfies \(\varphi\leq f\), we conclude
\begin{equation}\label{eq:ch23-bounded-obstacle-envelope}
 P_\theta(f)(x)=
 \begin{cases}
  0,& x\ne x_0,\\
  -1,& x=x_0 .
 \end{cases}
\end{equation}

Three conclusions may be read off \eqref{eq:ch23-bounded-obstacle-envelope}.
The raw envelope \(P_\theta(f)\) is not upper semicontinuous at \(x_0\): its
value there is \(-1\) while it equals \(0\) on every punctured neighborhood.
Its regularization is \(P_\theta(f)^*\equiv0\), which is \(\theta\)-psh but
is \emph{not} dominated by \(f\).  And there is no largest \(\theta\)-psh
function dominated by \(f\): such a function would have to be \(\geq
\chi_\varepsilon\) for every \(\varepsilon\), hence \(\geq0\) off \(x_0\),
hence \(\equiv 0\) by upper semicontinuity, contradicting
\(\varphi(x_0)\leq-1\).  Thus for bounded obstacles the star in
\Cref{def:ch23-envelope} is indispensable, and the phrase ``the largest
\(\theta\)-psh function dominated by \(f\)'' is simply not available.  The
mechanism is exactly that of \Cref{ex:sup-not-usc}: a supremum of upper
semicontinuous functions is upper semicontinuous only after regularization.
\end{example}

\subsection{Choquet selection on one dual complex}
\label{sec:ch23-choquet-on-delta}

The upper-envelope theorem was proved with nets because the index set of a
family of psh functions, and the system of models, are both possibly
uncountable.  On one fixed dual complex the situation is completely
different, and it is worth recording exactly why.

\begin{proposition}
\label{prop:ch23-skeleton-choquet}
Let \(\Delta=\Delta_{\mathcal X}\) be the dual complex of an SNC model, and
let \(\mathcal G\) be a nonempty family of functions
\(\Delta\to\Rbar\).  Then there is a countable subfamily
\(\{g_j\}_{j\geq1}\subset\mathcal G\) with
\begin{equation}\label{eq:ch23-choquet-on-delta}
 \Bigl(\sup_{g\in\mathcal G}g\Bigr)^{*}
 =\Bigl(\sup_{j\geq1}g_j\Bigr)^{*} ,
\end{equation}
the regularizations being taken in \(\Delta\); and if \(\mathcal G\) is
stable under finite maxima, the sequence may be taken increasing.
\end{proposition}

\begin{proof}
The dual complex is a finite union of simplices glued along faces, hence a
compact metrizable space, and in particular second countable.  Apply
\Cref{lem:choquet-selection}.
\end{proof}

\begin{remark}[The two countability regimes]
\label{rem:ch23-two-countability-regimes}
\Cref{prop:ch23-skeleton-choquet} is available on each \(\Delta\) and is not
available on \(X^{\An}\): a Berkovich analytification of positive dimension
over a discretely valued field is compact but not metrizable, and the
system of SNC models has no cofinal sequence, by
\Cref{rem:cofinality-scope} and \Cref{ex:no-cofinal-sequence}.  The correct
formulation of the difference is
\Cref{def:countably-regularizing}: what one would need on \(X^{\An}\) is a
countable regularizing family of open sets, and
\Cref{lem:regularizing-criterion} shows that this amounts to a countable
neighborhood basis at every point, which fails at every point of a skeleton
of positive dimension.  \Cref{ex:choquet-fails-uncountable} shows that the
conclusion, not merely the proof, breaks down without it.  The selection
principle of the next subsection is therefore stated, and can only be
stated, on one fixed dual complex.
\end{remark}

\subsection{A countable selection on one skeleton}
\label{sec:ch23-fixed-skeleton-selection}

The global theory requires nets, but one finite dual complex is metrizable.
The following selection principle is useful when an integral is supported
on a fixed skeleton.

\begin{proposition}\label{prop:ch23-fixed-skeleton-selection}
Let \(\mathcal F\subset\PSH(X,\theta)\) be nonempty and uniformly bounded
above, and fix a determining dual complex \(\Delta\).  Then
\begin{equation}\label{eq:ch23-fixed-skeleton-sup}
                    u:=\sup_{\psi\in\mathcal F}\psi|_\Delta
\end{equation}
is continuous.  There is an increasing sequence \(u_j\), each the
restriction of a finite maximum of members of \(\mathcal F\), such that
\(u_j\to u\) uniformly on \(\Delta\).
\end{proposition}

\begin{proof}
Write \(\Delta=\Delta_{\mathcal X}\), fix \(\psi_0\in\mathcal F\), and let
\(C_0\in\mathbf R\) be such that \(\psi\leq C_0\) on \(X^{\An}\) for every
\(\psi\in\mathcal F\).

\emph{Step 1: the exact skeletal inputs and a common global modulus.}
We use the following two earlier BFJ statements, specialized to the fixed
determining SNC model \(\mathcal X\).  First, clause~(1) of
\Cref{thm:bfj-skeletal-characterization} says that, for every
\(\psi\in\PSH(X,\theta)\), the restriction \(\psi|_\Delta\) is a finite
real-valued continuous function.  Second, clause~(1) of
\Cref{thm:bfj-skeletal-equicontinuity} says the following more precise
uniform statement.  Fix, on each face \(\sigma\) of \(\Delta\), a norm
\(|\cdot|_\sigma\) compatible with its integral affine lattice.  There is a
constant \(L=L(\mathcal X,\theta)>0\), independent of \(\psi\), such that
\begin{equation}\label{eq:ch23-selection-facewise}
 |\psi(x)-\psi(y)|\leq L|x-y|_\sigma,
 \qquad x,y\in\sigma,
 \quad \psi\in\PSH(X,\theta).
\end{equation}

We now pass from \eqref{eq:ch23-selection-facewise} to one metric on the
whole complex, including the topology check needed below.  For
\(x,y\in\Delta\), define
\begin{equation}\label{eq:ch23-selection-path-metric}
 d(x,y):=\inf\left\{
   \sum_{i=1}^{N}|x_i-x_{i-1}|_{\sigma_i}:
   x_0=x,\ x_N=y,\ [x_{i-1},x_i]\subset\sigma_i
 \right\}.
\end{equation}
In each summand we take \(\sigma_i\) to be the smallest face containing the
segment.  The set of chains in \eqref{eq:ch23-selection-path-metric} is
nonempty: \Cref{lem:dual-complex-simplicial} says in this setting that
\(\Delta\) is a connected finite simplicial complex, and such a complex is
joined by chains of segments lying in single faces.  Reversing and
concatenating chains show, respectively, that \(d\) is symmetric and
satisfies the triangle inequality, while \(d(x,x)=0\).

To verify that \(d\) is compatible with the simplicial topology, choose a
geometric realization of \(\Delta\) in a Euclidean space and denote its
Euclidean distance by \(\rho\).  Norm equivalence on the finitely many faces
gives constants \(a,b>0\) such that
\begin{equation}\label{eq:ch23-selection-norm-comparison}
 a\rho(p,q)\leq |p-q|_\sigma\leq b\rho(p,q)
 \qquad(p,q\in\sigma)
\end{equation}
for every face \(\sigma\).  The triangle inequality for \(\rho\), applied to
every admissible chain, yields \(d(x,y)\geq a\rho(x,y)\); hence \(d\) is
nondegenerate.  Conversely, for fixed \(x\), the union of the faces not
containing \(x\) is compact and misses \(x\).  Thus every sufficiently
\(\rho\)-close point \(y\) lies with \(x\) in one face \(\sigma\), and the
one-segment chain together with \eqref{eq:ch23-selection-norm-comparison}
gives \(d(x,y)\leq b\rho(x,y)\).  Therefore \(d\) induces precisely the
simplicial topology; in particular \((\Delta,d)\) is compact.

Apply \eqref{eq:ch23-selection-facewise} to every segment of a chain, sum,
and take the infimum in \eqref{eq:ch23-selection-path-metric}.  We obtain
\begin{equation}\label{eq:ch23-selection-modulus}
 |\psi(x)-\psi(y)|\leq L\,d(x,y),
 \qquad x,y\in\Delta,
 \quad \psi\in\PSH(X,\theta).
\end{equation}
We write \(\omega(t)=Lt\).  If \(f_1,\ldots,f_r\) satisfy this modulus, then
\[
 |\max_i f_i(x)-\max_i f_i(y)|
 \leq\max_i|f_i(x)-f_i(y)|
 \leq\omega(d(x,y)).
\]
Thus every finite maximum of restrictions of members of \(\mathcal F\) has
the same modulus.  No psh stability of that maximum is needed here.

\emph{Step 2: the supremum is finite and continuous.}
For \(x\in\Delta\) we have
\(\psi_0(x)\leq u(x)\leq C_0\).  The first exact skeletal input recorded in
Step~1 makes \(\psi_0(x)\) finite, so \(u\) is real valued.  For
\(x,y\in\Delta\) and \(\psi\in\mathcal F\),
\eqref{eq:ch23-selection-modulus} gives
\(\psi(x)\leq\psi(y)+\omega(d(x,y))\leq u(y)+\omega(d(x,y))\); taking the
supremum over \(\psi\) gives \(u(x)\leq u(y)+\omega(d(x,y))\), and
exchanging \(x\) and \(y\) gives the reverse.  Thus \(u\) itself satisfies
\eqref{eq:ch23-selection-modulus} and is in particular continuous.

\emph{Step 3: the dense-set selection.}
The dual complex is a finite simplicial complex, hence a compact metrizable,
and therefore separable, space; choose a countable dense subset
\((x_k)_{k\geq1}\subset\Delta\).  For every pair \(k,m\geq1\) use the
definition of the pointwise supremum \(u(x_k)\), which is a finite real
number by Step~2, to select \(\psi_{k,m}\in\mathcal F\) with
\begin{equation}\label{eq:ch23-dense-selection}
                     \psi_{k,m}(x_k)>u(x_k)-m^{-1}.
\end{equation}
Fix a bijection \(j\mapsto(k_j,m_j)\) from \(\mathbf N_{>0}\) to
\(\mathbf N_{>0}^2\), put \(\psi^{(j)}:=\psi_{k_j,m_j}\), and set
\begin{equation}\label{eq:ch23-selection-finite-maxima}
 u_j:=\bigl(\max\{\psi^{(1)},\ldots,\psi^{(j)}\}\bigr)\big|_\Delta,
 \qquad j\geq1 .
\end{equation}
Each \(u_j\) is by construction the restriction to \(\Delta\) of a finite
maximum of members of \(\mathcal F\); the sequence \((u_j)\) is increasing,
and \(u_j\leq u\) for every \(j\) because each \(\psi^{(i)}\leq u\) on
\(\Delta\).

\emph{Step 4: the limit is \(u\).}
Let \(w:=\lim_ju_j=\sup_ju_j\), a pointwise increasing limit, so
\(w\leq u\).  Each \(u_j\) satisfies \eqref{eq:ch23-selection-modulus} by
Step~1, hence so does \(w\), by the argument of Step~2 applied to the
countable family \(\{u_j\}\); in particular \(w\) is continuous.  At a point
\(x_k\) of the dense set, \eqref{eq:ch23-dense-selection} gives
\(u_j(x_k)>u(x_k)-m^{-1}\) as soon as \(j\) is large enough that
\(\psi_{k,m}\) occurs among \(\psi^{(1)},\ldots,\psi^{(j)}\); letting first
\(j\to\infty\) and then \(m\to\infty\) yields \(w(x_k)=u(x_k)\).  Two
continuous functions on \(\Delta\) that agree on a dense subset agree
everywhere, so \(w=u\).

\emph{Step 5: uniformity.}
We include the sequential Dini argument in the form needed here.  Fix
\(\varepsilon>0\) and set
\[
 V_j:=\{x\in\Delta:u_j(x)>u(x)-\varepsilon\}.
\]
The sets \(V_j\) are open because \(u_j-u\) is continuous, they increase
with \(j\), and they cover \(\Delta\) because \(u_j(x)\nearrow u(x)\) at
every point.  Compactness gives a finite subcover; since the sets are nested,
one of them, say \(V_N\), already equals \(\Delta\).  Consequently
\(0\leq u-u_j<\varepsilon\) on \(\Delta\) for every \(j\geq N\).  Hence
\(\|u-u_j\|_\infty\to0\), as asserted.
\end{proof}

\begin{remark}[Exactly what has and has not been selected]
\label{rem:ch23-selection-scope}
The proposition makes no assertion that the chosen sequence works on all
skeleta simultaneously.  Such an assertion would require a countable
cofinal model system and is not used.  We spell out the point, because it is
the single most tempting false step in this part of the theory.

The sequence \((u_j)\) produced in the proof depends on \(\Delta\) through
three separate choices: the constant \(L\) of
\eqref{eq:ch23-selection-modulus}, which comes from
\Cref{thm:bfj-skeletal-equicontinuity} and blows up as the model grows; the
countable dense subset \((x_k)\); and the selection
\eqref{eq:ch23-dense-selection} itself.  Nothing in the construction is
uniform in \(\Delta\).  Passing to a larger skeleton \(\Delta'\supset\Delta\)
one obtains a different sequence, and there is in general no single sequence
in \(\mathcal F\) whose finite maxima converge uniformly on every dual
complex: a sequence of \(\theta\)-psh functions determines, and is
determined by, countably many skeletal restrictions, whereas the system of
SNC models has no countable cofinal subsystem
(\Cref{rem:cofinality-scope}, \Cref{ex:no-cofinal-sequence}).

What \emph{is} true globally is \Cref{thm:ch23-upper-envelope}, which is a
statement about nets and whose conclusion involves a regularization.  The
present proposition is the sequential statement available after one skeleton
has been fixed, and it is used only in that form: for instance when a measure
is known to be supported on \(\Delta\), so that a single skeleton carries all
the information, as in \Cref{thm:bfj-supported-regularity}.
\end{remark}

\begin{remark}[Relation to Choquet's lemma]
\label{rem:ch23-selection-vs-choquet}
\Cref{prop:ch23-fixed-skeleton-selection} and
\Cref{prop:ch23-skeleton-choquet} are close but not identical, and the
difference matters.  Choquet's lemma
\Cref{lem:choquet-selection} produces a countable subfamily with the same
\emph{regularized} supremum, and it says nothing about the mode of
convergence.  The proposition above produces a countable subfamily whose
finite maxima converge \emph{uniformly} to the unregularized supremum.  The
extra strength comes entirely from the equicontinuity estimate
\eqref{eq:ch23-selection-modulus}: it is what makes the supremum continuous,
so that no star is needed, and it is what upgrades pointwise convergence on
a dense set to pointwise convergence everywhere.  Without a common modulus
one has only Choquet's conclusion, and
\Cref{ex:choquet-pointwise-fails} shows that pointwise equality of the
suprema can genuinely fail for the selected subfamily.
\end{remark}

\section{Continuity of continuous-obstacle envelopes}
\label{sec:ch23-continuity}

Being psh does not yet make \(P_\theta(f)\) continuous away from skeleta.
Everything proved so far controls \(P_\theta(f)\) on each dual complex, where
equicontinuity is available, and reconstructs it elsewhere as a decreasing
limit of retractions; a decreasing limit of continuous functions is upper
semicontinuous and nothing more.  The next theorem is the central algebraic
input of this chapter, and the first statement whose proof leaves the
category of skeleta.

\begin{center}
\fbox{\begin{minipage}{0.91\textwidth}
\textbf{Hypothesis signature of \Cref{thm:bfj-envelope-regularity}.}
\(K\) complete discretely valued, residue characteristic \(0\); \(X\) smooth
projective; \(\theta\) a closed model \((1,1)\)-form with \(\{\theta\}\)
ample; \(f\) continuous.  Each of these is used: discreteness in the model
theory of \Cref{chap:bfj-model-forms}, smoothness and residue characteristic
zero for the existence of log resolutions and of a cofinal system of SNC
models, ampleness in \Cref{lem:ch23-envelope-finiteness} and in the
rationalization step, and continuity of \(f\) in
\Cref{lem:ch23-retraction-uniform} and in the density argument.  What happens
when one of them is dropped is the subject of \Cref{sec:ch23-later}.
\end{minipage}}
\end{center}

\begin{theorem}\label{thm:bfj-envelope-regularity}
For every \(f\in C^0(X^{\An})\), the envelope \(P_\theta(f)\) is a uniform
limit on \(X^{\An}\) of \(\theta\)-psh model functions.  In particular,
\begin{equation}\label{eq:ch23-envelope-continuity}
                         P_\theta(f)\in C^0(X^{\An}).
\end{equation}
\end{theorem}

\begin{proof}
The proof has two independent ingredients, both established in full below.
\Cref{thm:ch23-base-ideal-approximation} proves uniform approximation for
an integral model reference and the zero obstacle.
\Cref{lem:ch23-reduction-to-base-ideals} then deduces the assertion for a
real model form and an arbitrary continuous obstacle by reference change,
rational approximation of the class, and a convex curvature correction.
Applying the theorem first and the lemma second produces, for every
prescribed error, a \(\theta\)-psh model function within that error of
\(P_\theta(f)\) on all of \(X^{\An}\).  Choosing errors tending to zero gives
uniform approximation.  Model functions are continuous, so their uniform
limit is continuous, proving \eqref{eq:ch23-envelope-continuity}.

There is no circularity: the proofs of the reduction lemma and the
base-ideal theorem use only the formal envelope calculus, model consistency,
and the multiplier-ideal package stated before the base-ideal theorem; they
do not invoke the present result.
\end{proof}

This is \cite[Theorem~8.3]{BFJ16}.  Its proof is not a consequence of
\Cref{thm:ch23-upper-envelope} or compactness.  We give the full reduction
and the numerical multiplier-ideal argument.  The exact algebraic input is
the uniform global-generation theorem stated before
\Cref{thm:ch23-base-ideal-approximation}.

Before the proof we record the two standard reformulations, both of which are
used constantly later and neither of which is a triviality.

\begin{corollary}
\label{cor:ch23-continuity-criterion}
Let \(f\in C^0(X^{\An})\) with \(\mathcal A_\theta(f)\ne\varnothing\), and
put
\begin{equation}\label{eq:ch23-model-competitor-family}
 \mathcal G_\theta(f):=\{\chi\in\mathcal D(X)_{\mathbf R}\cap\PSH(X,\theta):
   \chi\leq f\text{ on }X^{\An}\},
\end{equation}
so that \(P^{\mathcal D}_\theta(f)=\sup\mathcal G_\theta(f)\).  The
following are equivalent:
\begin{enumerate}
\item \(P^{\mathcal D}_\theta(f)\) is continuous on \(X^{\An}\);
\item \(P^{\mathcal D}_\theta(f)\) is a uniform limit of members of
\(\mathcal G_\theta(f)\);
\item \(P^{\mathcal D}_\theta(f)\) is a uniform limit of \(\theta\)-psh
model functions.
\end{enumerate}
If these hold, then \(P^{\mathcal D}_\theta(f)=P_\theta(f)\) on \(X^{\An}\).
\end{corollary}

\begin{proof}
\((2)\Rightarrow(3)\) is trivial, and \((3)\Rightarrow(1)\) holds because a
uniform limit of continuous functions is continuous and model functions are
continuous.

\((1)\Rightarrow(2)\).  The family \(\mathcal G_\theta(f)\) is nonempty by
\Cref{lem:ch23-envelope-finiteness} and stable under finite maxima by
\Cref{prop:ch20-model-max}, hence upward directed for the pointwise order.
It is therefore an increasing net of continuous functions with pointwise
supremum \(P^{\mathcal D}_\theta(f)\), which is continuous by hypothesis.
Apply \Cref{prop:dini-nets} to the decreasing net of continuous nonnegative
functions \(P^{\mathcal D}_\theta(f)-\chi\), \(\chi\in\mathcal G_\theta(f)\),
on the compact space \(X^{\An}\): its pointwise infimum is \(0\), so the
convergence is uniform.

For the final assertion, assume (1).  Each \(\chi\in\mathcal G_\theta(f)\)
satisfies \eqref{eq:ch21-retraction-inequality}, so
\(P^{\mathcal D}_\theta(f)\leq P^{\mathcal D}_\theta(f)\circ p_{\mathcal X}\)
for every determining SNC model, and
\eqref{eq:ch21-retraction-composition} makes the net
\(\bigl(P^{\mathcal D}_\theta(f)\circ p_{\mathcal X}\bigr)_{\mathcal X}\)
decreasing.  Since \(p_{\mathcal X}(x)\to x\) by
\Cref{thm:skeleton-inverse-limit} and \(P^{\mathcal D}_\theta(f)\) is
continuous, that net converges pointwise to
\(P^{\mathcal D}_\theta(f)\), whence
\(\inf_{\mathcal X}P^{\mathcal D}_\theta(f)(p_{\mathcal X}(x))
=P^{\mathcal D}_\theta(f)(x)\).  On the other hand \(P_\theta(f)\) is
\(\theta\)-psh by \Cref{cor:ch23-continuous-obstacle-psh}, so
\eqref{eq:ch21-skeletal-reconstruction} gives
\[
 P_\theta(f)(x)
 =\inf_{\mathcal X}P_\theta(f)\bigl(p_{\mathcal X}(x)\bigr)
 =\inf_{\mathcal X}P^{\mathcal D}_\theta(f)\bigl(p_{\mathcal X}(x)\bigr)
 =P^{\mathcal D}_\theta(f)(x),
\]
the middle equality by \Cref{prop:ch23-model-vs-psh-envelope}, since every
point \(p_{\mathcal X}(x)\) is quasimonomial.
\end{proof}

\begin{corollary}
\label{cor:ch23-envelope-by-model-functions}
For every \(f\in C^0(X^{\An})\),
\begin{equation}\label{eq:ch23-two-envelopes-equal}
 P^{\mathcal D}_\theta(f)=P_\theta(f)
 \qquad\text{on all of }X^{\An},
\end{equation}
and this common function is a uniform limit of members of
\(\mathcal G_\theta(f)\).
\end{corollary}

\begin{proof}
Let \(\varepsilon>0\).  By \Cref{thm:bfj-envelope-regularity} there is a
\(\theta\)-psh model function \(\chi\) with
\(\|\chi-P_\theta(f)\|_\infty<\varepsilon\).  Then
\(\chi-\varepsilon\leq P_\theta(f)\leq f\) on \(X^{\An}\), so
\(\chi-\varepsilon\in\mathcal G_\theta(f)\) and
\[
 P^{\mathcal D}_\theta(f)\geq\chi-\varepsilon>P_\theta(f)-2\varepsilon .
\]
Letting \(\varepsilon\downarrow0\) and combining with
\eqref{eq:ch23-model-envelope-below} gives
\eqref{eq:ch23-two-envelopes-equal}.  The second assertion is then
\Cref{cor:ch23-continuity-criterion}, whose hypothesis (1) now holds because
\(P^{\mathcal D}_\theta(f)=P_\theta(f)\) is continuous by
\eqref{eq:ch23-envelope-continuity}.
\end{proof}

\begin{corollary}
\label{cor:ch23-decreasing-regularization-pointer}
\Cref{thm:bfj-decreasing-regularization} holds: every
\(\varphi\in\PSH(X,\theta)\) is the pointwise limit of a decreasing net of
\(\theta\)-psh model functions.
\end{corollary}

\begin{proof}
This is exactly the deduction carried out in
\Cref{chap:bfj-compactness}: the order-theoretic argument there uses
\Cref{thm:bfj-envelope-regularity} as its only nonformal input, and that
theorem has now been proved.  We do not repeat the argument.
\end{proof}

\subsection{Reduction to a zero obstacle}

The reduction is most transparent when positivity is retained while the
reference class is rationalized.  We first isolate the density statement,
which is used twice.

\begin{lemma}
\label{lem:ch23-model-density}
The space \(\mathcal D(X)_{\mathbf R}\) of real model functions is uniformly
dense in \(C^0(X^{\An})\).
\end{lemma}

\begin{proof}
This is \Cref{prop:metrics-model-function-density}, which gives density of
the rational model functions \(\mathcal D(X)_{\mathbf Q}\); a fortiori the
larger space \(\mathcal D(X)_{\mathbf R}\) is dense.  The proof there is the
lattice form of the Stone--Weierstrass theorem: \(X^{\An}\) is compact
Hausdorff, model functions form a vector lattice
(\Cref{lem:metrics-model-lattice}) containing the constants
(\(\varphi_{c\mathcal X_s}=c\), by \eqref{eq:ch20-sign}) and separating the
points of \(X^{\An}\) (\Cref{lem:metrics-model-separation}).
\end{proof}

\begin{lemma}\label{lem:ch23-reduction-to-base-ideals}
Assume that the conclusion of
\Cref{thm:ch23-base-ideal-approximation} holds for every integral model
reference and the zero obstacle.  Then
\Cref{thm:bfj-envelope-regularity} holds for every continuous obstacle and
every real model form with ample generic class.
\end{lemma}

\begin{proof}
The proof is a chain of four reductions.  At each stage an error of size
\(O(\varepsilon)\) is committed, and the errors are added at the end; we
therefore fix \(\varepsilon>0\) once and for all and produce a
\(\theta\)-psh model function within \(C\varepsilon\) of \(P_\theta(f)\),
with \(C\) an absolute constant.

\emph{Reduction 1: to a positive reference form.}
By \Cref{prop:ch20-positive-representative} choose a model function \(v\)
such that \(\theta':=\theta+\ddc v\) is \(\mathcal X_0\)-positive for some
model \(\mathcal X_0\), that is, determined there by a relatively ample
class.  Part (4) of \Cref{prop:ch23-formal-envelope-calculus} gives
\(P_{\theta'}(f-v)=P_\theta(f)-v\), and
\eqref{eq:ch21-reference-class-example} identifies \(\theta'\)-psh model
functions with \(\theta\)-psh model functions minus \(v\).  So it is enough
to approximate \(P_{\theta'}(f-v)\) uniformly by \(\theta'\)-psh model
functions and then add \(v\).

\emph{Reduction 2: to a model obstacle.}
By \Cref{lem:ch23-model-density} choose a model function \(u\) with
\(\|u-(f-v)\|_\infty<\varepsilon\).  Both envelopes are finite by
\Cref{cor:ch23-continuous-finite}, so the contraction estimate
\eqref{eq:ch23-envelope-contraction} applies and gives
\begin{equation}\label{eq:ch23-obstacle-reduction}
       \|P_{\theta'}(u)-P_{\theta'}(f-v)\|_\infty<\varepsilon.
\end{equation}

\emph{Reduction 3: to a rational reference class.}
Pass to one model \(\mathcal X_1\geq\mathcal X_0\) determining both
\(\theta'\) and \(u\).  Rational classes are dense in the
finite-dimensional real vector space \(N^1(\mathcal X_1/R)\), so part (7) of
\Cref{prop:ch23-formal-envelope-calculus} allows us to choose a rational
class \(\theta_\varepsilon\in N^1(\mathcal X_1/R)_{\mathbf Q}\), as close to
\(\theta'\) as we please, with
\begin{equation}\label{eq:ch23-class-reduction}
 \|P_{\theta_\varepsilon}(u)-P_{\theta'}(u)\|_\infty<\varepsilon .
\end{equation}
We shall impose one further smallness condition on
\(\|\theta_\varepsilon-\theta'\|\) in Reduction~4; both are conditions of the
form ``sufficiently close'', so they can be met simultaneously.  Shrinking
if necessary, we also arrange that \(\theta_\varepsilon\) is
\(\mathcal X_1\)-positive, which is possible because ampleness is an open
condition.

\emph{Reduction 4: to an integral reference and the zero obstacle.}
Put \(\alpha_\varepsilon:=\theta_\varepsilon+\ddc u\), a rational class on
\(\mathcal X_1\); by construction it is relatively ample with ample generic
class.  Choose \(q\in\mathbf Z_{>0}\) with \(q\alpha_\varepsilon\) integral,
represented by a line bundle \(\mathcal M\) on \(\mathcal X_1\) whose
generic restriction is ample.  Homogeneity
\eqref{eq:ch23-homogeneity} gives
\(P_{q\alpha_\varepsilon}(0)=q\,P_{\alpha_\varepsilon}(0)\).  The assumed
conclusion of \Cref{thm:ch23-base-ideal-approximation}, applied to
\(\mathcal M\), supplies a \(q\alpha_\varepsilon\)-psh model function within
\(q\varepsilon\) of \(P_{q\alpha_\varepsilon}(0)\) in the uniform norm;
dividing by \(q\) gives an \(\alpha_\varepsilon\)-psh model function within
\(\varepsilon\) of \(P_{\alpha_\varepsilon}(0)\).  Adding \(u\) and using
reference change once more, we obtain a \(\theta_\varepsilon\)-psh model
function \(g_\varepsilon\) with
\begin{equation}\label{eq:ch23-rational-model-approximant}
       \|g_\varepsilon-P_{\theta_\varepsilon}(u)\|_\infty<\varepsilon.
\end{equation}

\emph{Correcting the curvature.}
The function \(g_\varepsilon\) is \(\theta_\varepsilon\)-psh, not
\(\theta'\)-psh, and \(\theta_\varepsilon\ne\theta'\).  The correction is a
convexity trick.  Choose \(t_\varepsilon\in(0,1)\) with
\(t_\varepsilon\downarrow0\) so slowly that
\(\|\theta_\varepsilon-\theta'\|/t_\varepsilon\to0\); this is the further
smallness condition promised in Reduction~3.  Set
\begin{equation}\label{eq:ch23-bracketing-form}
       \eta_\varepsilon
       :=\frac{\theta'-(1-t_\varepsilon)\theta_\varepsilon}
                    {t_\varepsilon}
        =\theta_\varepsilon
         +\frac{\theta'-\theta_\varepsilon}{t_\varepsilon},
\end{equation}
so that by construction
\begin{equation}\label{eq:ch23-bracketing-identity}
 \theta'=(1-t_\varepsilon)\theta_\varepsilon
        +t_\varepsilon\eta_\varepsilon .
\end{equation}
The second expression in \eqref{eq:ch23-bracketing-form} shows
\(\eta_\varepsilon\to\theta'\) in the finite-dimensional space
\(N^1(\mathcal X_1/R)\), because the chosen ratio tends to \(0\).  Since
\(\theta'\) is \(\mathcal X_1\)-positive and ampleness is open,
\(\eta_\varepsilon\) is relatively ample for all small \(\varepsilon\); in
particular the constant \(0\) is \(\eta_\varepsilon\)-psh.  Now
\eqref{eq:ch23-bracketing-identity} and the convexity of the semipositive
cone (\Cref{prop:ch20-semipositive-well-defined}) show that
\[
 \theta'+\ddc\bigl((1-t_\varepsilon)g_\varepsilon\bigr)
 =(1-t_\varepsilon)\bigl(\theta_\varepsilon+\ddc g_\varepsilon\bigr)
  +t_\varepsilon\bigl(\eta_\varepsilon+\ddc 0\bigr)
\]
is a convex combination of two semipositive classes, hence semipositive.
Therefore
\[
                   h_\varepsilon
                   :=(1-t_\varepsilon)g_\varepsilon
                    +t_\varepsilon\cdot 0
                   =(1-t_\varepsilon)g_\varepsilon
\]
is a \(\theta'\)-psh model function.

\emph{Assembling the errors.}
By \eqref{eq:ch23-obstacle-reduction},
\eqref{eq:ch23-class-reduction} and
\eqref{eq:ch23-rational-model-approximant},
\[
 \|g_\varepsilon-P_{\theta'}(f-v)\|_\infty<3\varepsilon,
\]
so the family \((g_\varepsilon)\) is uniformly bounded, say by \(B\),
independently of \(\varepsilon\), because \(P_{\theta'}(f-v)\) is a fixed
bounded function by \Cref{cor:ch23-continuous-finite}.  Hence
\(\|h_\varepsilon-g_\varepsilon\|_\infty
=t_\varepsilon\|g_\varepsilon\|_\infty\leq t_\varepsilon B\to0\), and
\[
 \|h_\varepsilon-P_{\theta'}(f-v)\|_\infty
 <3\varepsilon+t_\varepsilon B .
\]
Adding \(v\) produces \(\theta\)-psh model functions converging uniformly to
\(P_\theta(f)\), as required.
\end{proof}

\begin{remark}[Why the curvature correction is needed at all]
\label{rem:ch23-curvature-correction}
The natural attempt is to approximate \(\theta'\) by rational classes and to
declare the resulting model functions \(\theta'\)-psh.  That is false: a
\(\theta_\varepsilon\)-psh function is \(\theta'\)-psh only if
\(\theta'-\theta_\varepsilon\) happens to be semipositive, and a small
perturbation of a nef class in an arbitrary direction is not nef.  The
device in \eqref{eq:ch23-bracketing-form} replaces ``perturb and hope'' by an
exact convex identity: \(\theta'\) is written as a convex combination of
\(\theta_\varepsilon\) and a class \(\eta_\varepsilon\) which is genuinely
ample, and the corresponding convex combination of potentials is
\(\theta'\)-psh by an exact cone statement.  The cost is the factor
\((1-t_\varepsilon)\), which is harmless precisely because the approximants
are uniformly bounded.  The same manoeuvre is used, in the same place, by
\cite[proof of Theorem~8.3]{BFJ16} and \cite[Lemma~8.4]{GJKM19}.
\end{remark}

Thus it remains to treat
\begin{equation}\label{eq:ch23-integral-reference}
             \theta=c_1(\mathcal L,\|\cdot\|_{\mathcal L}),
\end{equation}
where \(\mathcal L\) is a line bundle on an SNC model \(\mathcal X\) and
\(L=\mathcal L|_X\) is ample.

\subsection{Base ideals}
\label{sec:ch23-base-ideals}

For \(m\gg1\), the generic bundle \(mL\) is globally generated.  Let
\(\mathfrak a_m\subset\mathcal O_{\mathcal X}\) be the vertical base ideal,
the image of
\begin{equation}\label{eq:ch23-base-ideal}
 H^0(\mathcal X,m\mathcal L)
       \otimes\mathcal O_{\mathcal X}(-m\mathcal L)
 \longrightarrow\mathcal O_{\mathcal X},
\end{equation}
and set
\begin{equation}\label{eq:ch23-base-potential}
                         \varphi_m
                         :=\frac1m\log|\mathfrak a_m|.
\end{equation}

Three remarks on \eqref{eq:ch23-base-ideal} before we proceed.  The ideal
\(\mathfrak a_m\) is the ideal of the base locus of \(m\mathcal L\) on the
model, computed with all its scheme structure.  It is \emph{vertical} for
\(m\gg1\), that is, its zero locus is contained in \(\mathcal X_s\): the
generic fibre of \(\mathfrak a_m\) is the base ideal of \(mL\), which is the
unit ideal once \(mL\) is globally generated.  Finally,
\(\log|\mathfrak a_m|\leq0\) because \(\mathfrak a_m\subset\cO_{\mathcal X}\),
by the sign convention \eqref{eq:ch20-sign}; hence \(\varphi_m\leq0\).

\subsection{The imported algebraic input}
\label{sec:ch23-imported-input}

Everything so far has been elementary.  The following theorem is the one
place in the chapter where a genuinely deep algebro-geometric result enters,
and we state it in the exact form in which it is used, with the exact
reference.  It is \emph{not} reproved here; its ingredients are reviewed in
\Cref{sec:appC-asymptotic-multiplier} and
\Cref{sec:appC-global-generation}.

\begin{theorem}
\label{thm:ch23-uniform-generation-input}
Let \(\mathcal X\) be a regular projective \(R\)-model of \(X\), with
\(K\) of residue characteristic zero.  Then:
\begin{enumerate}
\item every nonzero graded sequence
\(\mathfrak a_\bullet=(\mathfrak a_m)_{m\geq1}\) of vertical ideals on
\(\mathcal X\) has asymptotic multiplier ideals
\(\mathcal J(\mathfrak a_\bullet^m)\), themselves vertical, satisfying
\begin{equation}\label{eq:ch23-input-containment}
 \mathfrak a_m\subset\mathcal J(\mathfrak a_\bullet^m)
 \qquad(m\geq1)
\end{equation}
and the subadditivity property
\begin{equation}\label{eq:ch23-input-subadditivity}
 \mathcal J(\mathfrak a_\bullet^{m\ell})
   \subset\mathcal J(\mathfrak a_\bullet^m)^{\ell}
 \qquad(m,\ell\geq1);
\end{equation}
\item there is an ample line bundle \(\mathcal B\) on \(\mathcal X\),
depending only on \(\mathcal X\), such that for every such
\(\mathfrak a_\bullet\) and every line bundle \(\mathcal M\) on
\(\mathcal X\) for which \(\mathcal M^{\otimes m}\otimes\mathfrak a_m\) is
globally generated at all sufficiently divisible \(m\), the sheaf
\begin{equation}\label{eq:ch23-input-uniform-generation}
       \mathcal B\otimes\mathcal M^{\otimes m}\otimes
       \mathcal J(\mathfrak a_\bullet^m)
\end{equation}
is globally generated for every \(m\geq1\).
\end{enumerate}
\end{theorem}

\begin{proof}[Foundational input and scope]
Part (1) is the model-theoretic version of
\Cref{prop:appC-asymptotic-multiplier-properties}: containment
\eqref{eq:ch23-input-containment} is
\Cref{prop:appC-asymptotic-multiplier-properties}\,(1), and
\eqref{eq:ch23-input-subadditivity} is
\Cref{thm:appC-multiplier-subadditivity} passed to the stabilized
asymptotic ideals as in
\eqref{eq:appC-asymptotic-subadditivity}.  Both use a log resolution of
\(\mathfrak a_q\) on \(\mathcal X\) and local vanishing, hence residue
characteristic zero; the subadditivity theorem of Demailly--Ein--Lazarsfeld
\cite[Theorem~2.3]{DEL00} additionally uses smoothness of the diagonal.
Part (2) is the model version of \Cref{thm:appC-uniform-generation-multiplier},
proved by relative Nadel vanishing on a log resolution followed by
Castelnuovo--Mumford regularity relative to \(\Spec R\); the essential point
is that the twist \(\mathcal B\) does not grow with \(m\).  Verticality of
\(\mathcal J(\mathfrak a_\bullet^m)\) follows from verticality of
\(\mathfrak a_\bullet\), since a log resolution of a vertical ideal is an
isomorphism over the generic fibre.
\end{proof}

\begin{remark}[Exactly what is imported]
\label{rem:ch23-what-is-imported}
Three properties, and no others: containment
\eqref{eq:ch23-input-containment}, subadditivity
\eqref{eq:ch23-input-subadditivity}, and uniform global generation
\eqref{eq:ch23-input-uniform-generation} with a twist independent of \(m\).
In particular we do \emph{not} use a valuative description of
\(\mathcal J(\mathfrak a_\bullet^m)\), nor Skoda division, nor any statement
about semivaluations with nonzero support; compare
\Cref{rem:appC-asymptotic-valuative-safeguard} and
\Cref{rem:appC-supported-points}.  The corresponding list for the positive
characteristic proof, in which \(\mathcal J\) is replaced by an asymptotic
test ideal, is tabulated in \Cref{sec:appC-envelope-parallel}; that the two
lists coincide item by item is the reason the argument transposes, and
\Cref{check:appC-characteristic-firewall} is the reason the transposition is
not an identity of sheaves.
\end{remark}

\begin{theorem}\label{thm:ch23-base-ideal-approximation}
In the setting of \eqref{eq:ch23-integral-reference}, every sufficiently
large
\(\varphi_m\) is a \(\theta\)-psh model function and
\begin{equation}\label{eq:ch23-base-ideal-convergence}
                         \varphi_m\longrightarrow P_\theta(0)
                         \quad\text{uniformly on }X^{\An}.
\end{equation}
\end{theorem}

\begin{proof}
This is the argument of \cite[Theorem~8.5]{BFJ16}, now written with every
deduction from the imported algebraic input explicit.  It has five steps.
Steps~1--2 are formal; Step~3 identifies the envelope on the quasimonomial
locus and uses no multiplier ideal; Step~4 is where
\Cref{thm:ch23-uniform-generation-input} enters and produces the
quantitative sandwich; Step~5 removes the restriction to quasimonomial
points and converts the sandwich into a rate.

\emph{Step 1: each \(\varphi_m\) is a \(\theta\)-psh model function.}
Fix \(m\) so large that \(mL\) is globally generated and
\(\mathfrak a_m\) is vertical.  By construction
\(m\mathcal L\otimes\mathfrak a_m\) is the image of
\(H^0(\mathcal X,m\mathcal L)\otimes\cO_{\mathcal X}\to m\mathcal L\), hence
is globally generated; a fortiori it is relatively globally generated over
\(\Spec R\).  \Cref{lem:ch20-generated-ideal-psh} therefore shows that
\(\log|\mathfrak a_m|\) is an \(m\theta\)-psh model function, where
\(m\theta=c_1(m\mathcal L,\|\cdot\|)\).  Multiplying a form and a potential
by \(m^{-1}>0\) preserves relative nefness, since the nef cone is a cone;
hence \(\varphi_m=m^{-1}\log|\mathfrak a_m|\) is \(\theta\)-psh.  As noted
after \eqref{eq:ch23-base-potential}, \(\varphi_m\leq0\), so
\(\varphi_m\in\mathcal A_\theta(0)\) and
\begin{equation}\label{eq:ch23-varphi-m-below-envelope}
 \varphi_m\leq P_\theta(0)\leq 0
 \qquad\text{on }X^{\An}.
\end{equation}

\emph{Step 2: the potentials are superadditive.}
Multiplying sections gives
\begin{equation}\label{eq:ch23-graded-base-ideals}
                \mathfrak a_m\mathfrak a_\ell
                \subset\mathfrak a_{m+\ell},
\end{equation}
because a product of a section of \(m\mathcal L\) and a section of
\(\ell\mathcal L\) is a section of \((m+\ell)\mathcal L\).  Since
\(\log|\mathfrak c\mathfrak d|=\log|\mathfrak c|+\log|\mathfrak d|\) and
\(\mathfrak c\subset\mathfrak d\) implies
\(\log|\mathfrak c|\leq\log|\mathfrak d|\), we get
\begin{equation}\label{eq:ch23-superadditive-potentials}
             m\varphi_m+\ell\varphi_\ell
             \leq(m+\ell)\varphi_{m+\ell}.
\end{equation}
Thus \((m\varphi_m(x))_m\) is a superadditive sequence of real numbers,
bounded above by \(0\); Fekete's lemma identifies its limit along
sufficiently large indices with \(\sup_m\varphi_m(x)\), pointwise on
\(X^{\An}\).

\emph{Step 3: identification of the envelope at quasimonomial points.}
We claim that
\begin{equation}\label{eq:ch23-envelope-base-sup}
                         P_\theta(0)(x)=\sup_m\varphi_m(x)
\end{equation}
at every quasimonomial \(x\).  In view of
\eqref{eq:ch23-varphi-m-below-envelope} only ``\(\leq\)'' needs proof.

\emph{Step 3a: a model competitor almost attaining the envelope at \(x\).}
Given \(\varepsilon>0\), we produce a \(\theta\)-psh model function
\(\varphi_D\) with
\begin{equation}\label{eq:ch23-model-competitor-at-point}
             \varphi_D\leq 0\text{ on }X^{\An},
             \qquad
             \varphi_D(x)\geq P_\theta(0)(x)-\varepsilon.
\end{equation}
Choose \(\psi\in\PSH(X,\theta)\) with \(\psi\leq0\) and
\(\psi(x)>P_\theta(0)(x)-\varepsilon\).  Let \(\mathcal Y\) be an SNC model,
determining \(\theta\), whose dual complex contains \(x\).  By clause (2) of
\Cref{def:ch21-bfj-psh} there is a \(\theta\)-psh model function
\(\varphi_D\) with
\[
 \left|\varphi_D-(\psi-\varepsilon)\right|<\varepsilon
 \quad\text{on }\Delta_{\mathcal Y};
\]
after enlarging \(\mathcal Y\) we may assume \(\varphi_D\) is determined on
\(\mathcal Y\).  On \(\Delta_{\mathcal Y}\) this gives
\(\varphi_D<\psi\leq0\) and
\(\varphi_D>\psi-2\varepsilon>P_\theta(0)(x)-3\varepsilon\) at \(x\).
Because \(\varphi_D\) is determined on \(\mathcal Y\) we have
\(\varphi_D=\varphi_D\circ p_{\mathcal Y}\)
(\Cref{lem:ch21-model-skeletal-properties}), whence for every
\(x'\in X^{\An}\)
\[
       \varphi_D(x')
       =\varphi_D(p_{\mathcal Y}(x'))
       \leq\psi(p_{\mathcal Y}(x'))\leq0 ,
\]
the last inequality because \(\psi\leq0\) everywhere.  So \(\varphi_D\leq0\)
on all of \(X^{\An}\).  Replacing \(3\varepsilon\) by \(\varepsilon\) at the
outset gives \eqref{eq:ch23-model-competitor-at-point}.  Note the shape of
this argument: a global inequality was obtained from a skeletal one by
composing with the retraction, and this is possible only because the
obstacle, here the constant \(0\), is itself constant on the fibres of
\(p_{\mathcal Y}\).  For a general continuous obstacle the same step
requires \Cref{lem:ch23-retraction-uniform}, as in
\Cref{prop:ch23-model-vs-psh-envelope}.

\emph{Step 3b: from a model competitor to a section.}
Let \(\pi:\mathcal X'\to\mathcal X\) dominate both \(\mathcal X\) and
\(\mathcal Y\), and regard \(\varphi_D=\varphi_{D}\) as attached to a
vertical \(\mathbf Q\)-divisor \(D\) on \(\mathcal X'\).  By
\Cref{prop:ch20-positive-representative}, applied to the semipositive class
\(\pi^*\theta+[D]\), there is on a further model a vertical
\(\mathbf Q\)-divisor \(D'\) with
\begin{equation}\label{eq:ch23-small-ample-correction}
     -\varepsilon\leq\varphi_{D'}\leq0,
     \qquad
     \pi^*\mathcal L+D+D'\text{ ample},
\end{equation}
where we have kept the name \(\mathcal X'\) for that further model.  Both
\(D\) and \(D'\) are anti-effective: a vertical divisor is anti-effective
exactly when its model function is nonpositive at every divisorial point of
a determining model, by \eqref{eq:ch20-sign}, and
\(\varphi_D\leq0\), \(\varphi_{D'}\leq0\) there.

Choose \(m\) divisible enough that \(m(\pi^*\mathcal L+D+D')\) is an ample
\emph{integral} line bundle and is globally generated.  Anti-effectivity of
\(D+D'\) gives an inclusion of sheaves
\begin{equation}\label{eq:ch23-section-inclusion}
 \mathcal O_{\mathcal X'}
       \bigl(m\pi^*\mathcal L+m(D+D')\bigr)
 \subset
 \mathcal O_{\mathcal X'}(m\pi^*\mathcal L),
\end{equation}
whose source is globally generated.  Since \(\mathcal X\) is normal and
\(\pi\) is birational, \(\pi_*\cO_{\mathcal X'}=\cO_{\mathcal X}\) and the
projection formula identify
\(H^0(\mathcal X',m\pi^*\mathcal L)=H^0(\mathcal X,m\mathcal L)\); so every
section occurring in \eqref{eq:ch23-section-inclusion} is pulled back from
\(\mathcal X\), and the evaluation map \eqref{eq:ch23-base-ideal} defining
\(\mathfrak a_m\) yields
\begin{equation}\label{eq:ch23-generated-inside-base-ideal}
       \mathcal O_{\mathcal X'}(m(D+D'))
       \subset\mathfrak a_m\mathcal O_{\mathcal X'} .
\end{equation}
Applying \(\log|\cdot|\) to
\eqref{eq:ch23-generated-inside-base-ideal}, which preserves inclusions, and
using \(\log|\cO(m(D+D'))|=m(\varphi_D+\varphi_{D'})\) together with
\(\log|\mathfrak a_m\cO_{\mathcal X'}|=\log|\mathfrak a_m|=m\varphi_m\),
gives \(\varphi_D+\varphi_{D'}\leq\varphi_m\) on \(X^{\An}\).  Evaluating at
\(x\), and using \(\varphi_{D'}\geq-\varepsilon\) and
\eqref{eq:ch23-model-competitor-at-point} this is
\begin{equation}\label{eq:ch23-base-lower-bound}
 P_\theta(0)(x)
 \leq\varphi_D(x)+\varepsilon
 \leq\varphi_m(x)+2\varepsilon.
\end{equation}
As \(\varepsilon>0\) was arbitrary and \(m\) may be taken arbitrarily
divisible, \eqref{eq:ch23-base-lower-bound} together with
\eqref{eq:ch23-varphi-m-below-envelope} proves
\eqref{eq:ch23-envelope-base-sup}.

\emph{Step 4: the quantitative sandwich.}
Pointwise convergence on quasimonomial points is not yet uniform, and no
rate has been obtained.  This is what the multiplier ideals supply.
Introduce
\begin{equation}\label{eq:ch23-asymptotic-multiplier}
                         \mathfrak b_m
                         :=\mathcal J(\mathfrak a_\bullet^m),
\end{equation}
the asymptotic multiplier ideals of the graded sequence
\(\mathfrak a_\bullet\) of \eqref{eq:ch23-graded-base-ideals}, which exist
and are vertical by \Cref{thm:ch23-uniform-generation-input}\,(1).  Their
two properties used here are
\begin{equation}\label{eq:ch23-multiplier-properties}
 \mathfrak a_m\subset\mathfrak b_m,
 \qquad
 \mathfrak b_{m\ell}\subset\mathfrak b_m^\ell,
\end{equation}
namely \eqref{eq:ch23-input-containment} and
\eqref{eq:ch23-input-subadditivity}.

\emph{Step 4a: the multiplier potential dominates the envelope.}
Combining the two inclusions of \eqref{eq:ch23-multiplier-properties},
\[
       \mathfrak a_{m\ell}
       \subset\mathfrak b_{m\ell}
       \subset\mathfrak b_m^\ell
       \qquad(m,\ell\geq1).
\]
Inclusion of ideals reverses no signs for our convention:
\(\mathfrak c\subset\mathfrak d\) implies
\(\log|\mathfrak c|\leq\log|\mathfrak d|\), by
\Cref{cor:retraction-ideal-monotone} and the definition of
\(\log|\cdot|\) through a normalized blowup.  Hence
\(m\ell\,\varphi_{m\ell}=\log|\mathfrak a_{m\ell}|\leq
\ell\log|\mathfrak b_m|\); dividing by \(m\ell\) and taking the supremum
over \(\ell\), then using \eqref{eq:ch23-envelope-base-sup}, gives on the
quasimonomial locus
\begin{equation}\label{eq:ch23-multiplier-upper-control}
              \frac1m\log|\mathfrak b_m|
              \geq P_\theta(0).
\end{equation}

\emph{Step 4b: uniform global generation converts \(\mathfrak b_m\) back
into a base ideal.}
After raising the model we may assume \(\mathcal X\) regular, and we choose
an effective vertical \(\mathbf Q\)-divisor \(E\) such that
\(\mathcal A:=\mathcal L-E\) is ample.  This is possible: \(\mathcal L\) has
ample generic restriction, so \Cref{prop:ch20-positive-representative}
supplies a model function \(\varphi_D\), which we may normalize to be
\(\leq0\) by adding a multiple of \(\mathcal X_s\) without changing any
numerical class, such that \(\mathcal L+D\) is relatively ample; and
\(\varphi_D\leq0\) at all divisorial points of a determining model means
exactly that \(D\) is anti-effective, so \(E:=-D\) is effective and
\(\mathcal A=\mathcal L-E\).  Apply
\Cref{thm:ch23-uniform-generation-input}\,(2) with
\(\mathcal M=\mathcal L\): there is an ample \(\mathcal B\), independent of
\(m\), with \(\mathcal B\otimes m\mathcal L\otimes\mathfrak b_m\) globally
generated for all \(m\).  Choose \(m_0\) so large that
\(m_0\mathcal A-\mathcal B\) is integral and globally generated; since a
tensor product of two globally generated sheaves is globally generated, and
\[
 \cO_{\mathcal X}(m\mathcal L+m_0\mathcal A)\otimes\mathfrak b_m
 =\bigl(\mathcal B\otimes m\mathcal L\otimes\mathfrak b_m\bigr)
  \otimes\bigl(m_0\mathcal A-\mathcal B\bigr),
\]
we obtain one integer \(m_0\), independent of \(m\), such that
\begin{equation}\label{eq:ch23-uniform-global-generation}
 \mathcal O_{\mathcal X}(m\mathcal L+m_0\mathcal A)
       \otimes\mathfrak b_m
 \quad\text{is globally generated for every }m.
\end{equation}
Since \(m_0\mathcal A=m_0\mathcal L-m_0E\), multiplication by the canonical
section of the effective divisor \(m_0E\) injects the globally generated
sheaf \(\mathcal F:=\cO_{\mathcal X}(m\mathcal L+m_0\mathcal A)\otimes
\mathfrak b_m\) of \eqref{eq:ch23-uniform-global-generation} into
\((m+m_0)\mathcal L\).  Its global sections therefore give sections of
\((m+m_0)\mathcal L\), which generate the subsheaf
\(\mathcal F\subset(m+m_0)\mathcal L\); the ideal they cut out is
\[
 \mathcal F\otimes\cO_{\mathcal X}\bigl(-(m+m_0)\mathcal L\bigr)
 =\mathfrak b_m\otimes\cO_{\mathcal X}(-m_0E) .
\]
Since \(\mathfrak a_{m+m_0}\) is generated by \emph{all} sections of
\((m+m_0)\mathcal L\), this ideal is contained in
\(\mathfrak a_{m+m_0}\); after pulling back to a model on which both ideals
are principal, and using
\(\log|\cO(-m_0E)|=-m_0\varphi_E\), one obtains
\begin{equation}\label{eq:ch23-global-generation-comparison}
 \log|\mathfrak b_m|
       \leq\log|\mathfrak a_{m+m_0}|+m_0\varphi_E .
\end{equation}
The term \(m_0\varphi_E\) is the price of the twist; note that it does not
grow with \(m\), which is the whole point of the fixed twist in
\Cref{thm:ch23-uniform-generation-input}\,(2) and of
\Cref{thm:appC-uniform-generation-multiplier}.

\emph{Step 4c: the sandwich.}
Substituting \eqref{eq:ch23-global-generation-comparison} into
\eqref{eq:ch23-multiplier-upper-control} with \(m\) replaced by \(m-m_0\)
gives
\begin{equation}\label{eq:ch23-uniform-sandwich}
 \varphi_m
 \leq P_\theta(0)
 \leq
 \frac{m}{m-m_0}\varphi_m
 +\frac{m_0}{m-m_0}\varphi_E,
 \qquad m\gg1 .
\end{equation}

\emph{Step 5: from quasimonomial points to \(X^{\An}\), and the rate.}
Inequality \eqref{eq:ch23-uniform-sandwich} has so far been proved on the
quasimonomial locus.  Its left inequality is
\eqref{eq:ch23-varphi-m-below-envelope} and already holds globally.  For the
right inequality, call the model function on the right \(g_m\) and choose a
model \(\mathcal Y\) determining it.  For \(x\in X^{\An}\) and every
\(\mathcal Z\geq\mathcal Y\), the point \(p_{\mathcal Z}(x)\) is
quasimonomial, so
\[
       P_\theta(0)(p_{\mathcal Z}(x))
       \leq g_m(p_{\mathcal Z}(x))=g_m(x),
\]
using \(g_m=g_m\circ p_{\mathcal Z}\).  Letting \(\mathcal Z\) increase and
using the skeletal reconstruction
\eqref{eq:ch21-skeletal-reconstruction} for the \(\theta\)-psh function
\(P_\theta(0)\), which is legitimate because \(P_\theta(0)\in\PSH(X,\theta)\)
by \Cref{cor:ch23-continuous-obstacle-psh}, extends the inequality to \(x\).
Thus no continuity of \(P_\theta(0)\) has been assumed, and this is the
step that would be circular if it were replaced by an appeal to
\Cref{thm:bfj-envelope-regularity}.

It remains to bound the right-hand side of
\eqref{eq:ch23-uniform-sandwich}.  Choose \(r\) so that \(\mathfrak a_s\) is
vertical and \(sL\) globally generated for every \(s\geq r\), and write a
sufficiently large \(m\) as \(m=kr+s\) with \(r\leq s<2r\) and
\(k\geq0\).  Iterating \eqref{eq:ch23-superadditive-potentials} gives
\[
 \varphi_m
 \geq\frac{kr}{m}\varphi_r+\frac{s}{m}\varphi_s .
\]
Only the finitely many functions \(\varphi_s\), \(r\leq s<2r\), occur on the
right, and each is a bounded model function; together with
\(\varphi_m\leq0\) this bounds the family \((\varphi_m)_{m\gg1}\) uniformly
in the sup norm, say \(\|\varphi_m\|_\infty\leq B\).  The model function
\(\varphi_E\) is bounded as well.  Subtracting \(\varphi_m\) throughout
\eqref{eq:ch23-uniform-sandwich}, which is legitimate because all terms are
finite, gives
\begin{equation}\label{eq:ch23-sandwich-difference}
 0\leq P_\theta(0)-\varphi_m
 \leq\frac{m_0}{m-m_0}
          \bigl(\varphi_m+\varphi_E\bigr) .
\end{equation}
Taking suprema and using the uniform bound just obtained,
\begin{equation}\label{eq:ch23-envelope-rate}
 \lVert P_\theta(0)-\varphi_m\rVert_\infty
 \leq \frac{m_0}{m-m_0}
 \bigl(B+\lVert\varphi_E\rVert_\infty\bigr)
 =O(m^{-1}).
\end{equation}
This proves the uniform convergence in
\eqref{eq:ch23-base-ideal-convergence}, with an explicit rate.

The only imported steps were
\Cref{thm:ch23-uniform-generation-input}.  Its construction from log
resolutions, relative Kawamata--Viehweg vanishing, and
Castelnuovo--Mumford regularity is reviewed in
\Cref{sec:appC-asymptotic-multiplier} and
\Cref{sec:appC-global-generation}.
\end{proof}

\begin{corollary}
\label{cor:ch23-envelope-rate}
In the setting of \eqref{eq:ch23-integral-reference} there are constants
\(m_0\) and \(C\), depending only on \(\mathcal X\) and \(\mathcal L\), with
\begin{equation}\label{eq:ch23-explicit-rate}
 0\leq P_\theta(0)-\varphi_m\leq\frac{C}{m}
 \qquad\text{on }X^{\An},\ m>m_0 .
\end{equation}
\end{corollary}

\begin{proof}
This is \eqref{eq:ch23-envelope-rate} with
\(C=2m_0(B+\|\varphi_E\|_\infty)\), valid for \(m\geq2m_0\).
\end{proof}

\begin{remark}[Where each hypothesis was used in the proof]
\label{rem:ch23-base-ideal-hypotheses}
Ampleness of \(L\) was used in Step~1, to have \(mL\) globally generated and
\(\mathfrak a_m\) vertical, and in Step~4b, to produce \(E\) with
\(\mathcal L-E\) ample.  Normality of the models was used in Step~3b, in the
identification \(\pi_*\cO_{\mathcal X'}=\cO_{\mathcal X}\).  Residue
characteristic zero was used exactly once, inside
\Cref{thm:ch23-uniform-generation-input}, and nowhere else; the arguments of
Steps~1--3 and~5 are characteristic free.  Discreteness of the valuation is
used throughout the model theory of \Cref{chap:bfj-model-forms} on which the
whole chapter rests.  The rate \(O(m^{-1})\) in
\eqref{eq:ch23-explicit-rate} is not used later, but it is worth recording,
because it is the same \(O(m^{-1})\) drift in numerical class that appears
in the trivially valued envelope proof of
\Cref{chap:bj-smooth-characteristic}.
\end{remark}

Together with \Cref{cor:ch23-decreasing-regularization-pointer}, the order
argument in the proof of
\Cref{thm:bfj-decreasing-regularization} converts
\Cref{thm:bfj-envelope-regularity} into decreasing model regularization.

\section{Orthogonality is a later theorem}

Continuity of an envelope says that the obstacle projection stays in
\(C^0\).  Orthogonality says where its Monge--Amp\`ere measure lives.  The
latter cannot be deduced by raising one vertical coefficient: that
perturbation can immediately leave the nef cone.

After the Monge--Amp\`ere operator has been constructed in
\Cref{chap:bfj-ma}, one can formulate the orthogonality property as

\begin{equation}\label{eq:ch23-orthogonality-preview}
 \int_{X^{\An}}
       (f-P_\theta(f))
       \MA_\theta(P_\theta(f))=0,
 \qquad f\in C^0(X^{\An}).
\end{equation}

The integrand is nonnegative, so \eqref{eq:ch23-orthogonality-preview}
means that the measure is carried by the contact set.  We do not assert
\eqref{eq:ch23-orthogonality-preview} from the work of this chapter.
\Cref{thm:orthogonality-provenance} proves it through asymptotic
orthogonality and differentiability of non-Archimedean volume, under the
precise historical hypotheses recorded there.

The dependency chain is therefore
\begin{equation}\label{eq:ch23-dependency-chain}
 \begin{aligned}
  &\text{skeletal equicontinuity}
    \Longrightarrow \text{supremum stability},\\
  &\text{multiplier ideals}
    \Longrightarrow \text{envelope continuity},\\
  &\text{envelope continuity}
    \Longrightarrow \text{decreasing regularization},\\
  &\left.
    \begin{gathered}
      \text{bounded Monge--Amp\`ere theory},\\[-2pt]
      \text{energy differentiability},\\[-2pt]
      \text{asymptotic orthogonality}
    \end{gathered}
    \right\}
    \Longrightarrow \eqref{eq:ch23-orthogonality-preview}.
 \end{aligned}
\end{equation}

This separation also prevents circularity: orthogonality is used later to
differentiate energy after projection, but it is not used to prove that the
projection is continuous.

\section{Orthogonality at model level}
\label{sec:ch23-model-orthogonality}

The previous section explained what orthogonality is not: it is not a
formal consequence of maximality.  This section explains what it is in the
two situations in which the obstacle problem is finite dimensional and
linear, namely on a curve model and in the polyhedral convex category.  In
both cases orthogonality is a complementary-slackness statement and has a
complete elementary proof.  In both cases the proof visibly fails in higher
dimension, and the reason it fails is the content of
\Cref{rem:ch23-model-slackness-fails}.

\subsection{The model obstacle problem}
\label{sec:ch23-model-obstacle-problem}

Throughout this subsection \(\mathcal X\) is a fixed normal projective model
and all model functions are determined on \(\mathcal X\).  Restricting the
competitors in \eqref{eq:ch23-envelope} to that finite-dimensional space
gives the \emph{model-level envelope}.

\begin{definition}[Model-level envelope on a fixed model]
\label{def:ch23-fixed-model-envelope}
Let \(\theta\) be determined on \(\mathcal X\) and let
\(u\in\mathcal D(X)_{\mathbf R}\) be determined on \(\mathcal X\).  Put
\begin{equation}\label{eq:ch23-fixed-model-competitors}
 K_{\mathcal X}(\theta,u)
 :=\{\varphi\in\mathcal D(X)_{\mathbf R}\ \text{determined on }\mathcal X:
     \varphi\leq u,\ \theta+\ddc\varphi\ \text{semipositive}\}
\end{equation}
and, when this set has a largest element, denote it by
\(P^{\mathcal X}_\theta(u)\).
\end{definition}

By construction
\begin{equation}\label{eq:ch23-model-level-below}
 P^{\mathcal X}_\theta(u)\leq P^{\mathcal X'}_\theta(u)
 \leq P^{\mathcal D}_\theta(u)=P_\theta(u),
 \qquad \mathcal X'\geq\mathcal X,
\end{equation}
the last equality being
\Cref{cor:ch23-envelope-by-model-functions}; the model-level envelope
increases with the model, and in general only its limit is the true
envelope.  \Cref{ex:ch23-curve-envelope-numbers} exhibits a case in which it
is already attained at a finite level, and
\Cref{rem:ch23-model-dependence} explains why that is a one-dimensional
phenomenon.

\subsection{Curves: complementary slackness on the dual graph}
\label{sec:ch23-curve-orthogonality}

Let \(n=1\), so that \(X\) is a smooth projective curve, and let
\(\mathcal X\) be a regular model with
\begin{equation}\label{eq:ch23-curve-special-fibre}
 \mathcal X_s=\sum_{i\in I}b_iE_i,
 \qquad
 w_{ij}:=b_ib_j(E_i\cdot E_j)\ \ (i\ne j),
 \qquad
 c_i:=b_i(\theta\cdot E_i).
\end{equation}
The numbers \(w_{ij}=w_{ji}\) are nonnegative, and they are positive exactly
when \(E_i\) and \(E_j\) meet; the graph \(\Gamma\) on the vertex set \(I\)
with those edges is the dual graph, and it is connected because
\(\mathcal X_s\) is.  If \(\theta\) is semipositive on \(\mathcal X\), then
\(c_i\geq0\) and, by \eqref{eq:ch20-special-fibre-degree},
\begin{equation}\label{eq:ch23-curve-degree-sum}
 \sum_{i\in I}c_i=\deg_X\{\theta\} .
\end{equation}
A model function determined on \(\mathcal X\) is the same thing as a vector
\(v=(v_i)_{i\in I}\in\mathbf R^I\) of values at the divisorial points
\(e_i:=x_{E_i}\), via \(v_i=\varphi(e_i)\); the corresponding vertical
divisor is \(\sum_ib_iv_iE_i\).  For such a \(v\) put
\begin{equation}\label{eq:ch23-curve-residual-degree}
 A_i(v):=c_i+\sum_{j\ne i}w_{ij}(v_j-v_i),
 \qquad i\in I .
\end{equation}

\begin{lemma}
\label{lem:ch23-curve-residual-degree}
With the notation \eqref{eq:ch23-curve-special-fibre}, for every
\(v\in\mathbf R^I\) with associated model function \(\varphi_v\):
\begin{enumerate}
\item \(A_i(v)=b_i\bigl((\theta+\ddc\varphi_v)\cdot E_i\bigr)\);
\item \(\theta+\ddc\varphi_v\) is semipositive if and only if
\(A_i(v)\geq0\) for every \(i\in I\);
\item the model Chambert--Loir measure of \(\theta+\ddc\varphi_v\) is
\begin{equation}\label{eq:ch23-curve-model-measure}
 \MA_\theta(\varphi_v)=\sum_{i\in I}A_i(v)\,\delta_{e_i},
 \qquad
 \MA_\theta(\varphi_v)(X^{\An})=\deg_X\{\theta\} .
\end{equation}
\end{enumerate}
\end{lemma}

\begin{proof}
(1) The vertical divisor of \(\varphi_v\) is \(D=\sum_jb_jv_jE_j\), so
\eqref{eq:ch20-graph-laplacian} gives
\((D\cdot E_i)=\sum_{j\ne i}b_j(E_j\cdot E_i)(v_j-v_i)\); multiplying by
\(b_i\) and adding \(c_i=b_i(\theta\cdot E_i)\) gives (1).

(2) In dimension one the proper integral curves contained in
\(\mathcal X_s\) are exactly the components \(E_i\), because
\(\mathcal X_s\) is one dimensional.  Semipositivity therefore means
nonnegativity of \(((\theta+\ddc\varphi_v)\cdot E_i)\) for all \(i\), which
by (1) and \(b_i>0\) is \(A_i(v)\geq0\).  This is
\eqref{eq:ch20-curve-nef-test}.

(3) The model measure of a model metric is
\(\sum_ib_i\deg\bigl((\theta+\ddc\varphi_v)|_{E_i}\bigr)\delta_{e_i}\) by
\Cref{thm:cl-model-properties} and \Cref{prop:ch08-curve-cl}; by (1) this is
\eqref{eq:ch23-curve-model-measure}.  The mass statement is
\Cref{thm:cl-model-properties}; alternatively, summing
\eqref{eq:ch23-curve-residual-degree} over \(i\) and using
\(w_{ij}=w_{ji}\) makes the double sum cancel in pairs, leaving
\(\sum_ic_i=\deg_X\{\theta\}\) by \eqref{eq:ch23-curve-degree-sum}.
\end{proof}

The last computation is worth isolating, because it is the mechanism behind
everything in this subsection: \(v\mapsto(A_i(v))_i\) is an affine map whose
linear part is minus a weighted graph Laplacian, and a weighted graph
Laplacian has nonnegative off-diagonal entries and vanishing row sums.

\begin{theorem}
\label{thm:ch23-model-orthogonality-curve}
Let \(X\) be a smooth projective curve, let \(\mathcal X\) be a regular
model, let \(\theta\) be a semipositive model form determined on
\(\mathcal X\), and let \(u\) be a real model function determined on
\(\mathcal X\), with vertex values \(u_i:=u(e_i)\).  Then:
\begin{enumerate}
\item the set \(K_{\mathcal X}(\theta,u)\) of
\eqref{eq:ch23-fixed-model-competitors} is nonempty and has a largest
element \(P:=P^{\mathcal X}_\theta(u)\);
\item \emph{(complementary slackness)} for every \(i\in I\),
\begin{equation}\label{eq:ch23-complementary-slackness}
 A_i(P)\,\bigl(u_i-P_i\bigr)=0 ;
\end{equation}
\item consequently
\begin{equation}\label{eq:ch23-model-orthogonality}
 \int_{X^{\An}}\bigl(u-P\bigr)\,\MA_\theta(P)=0,
\end{equation}
that is, \(\MA_\theta(P)\) is carried by the contact set
\(\{u=P\}\).
\end{enumerate}
\end{theorem}

\begin{proof}
Throughout we identify a model function determined on \(\mathcal X\) with its
vector of vertex values, and we use \Cref{lem:ch23-curve-residual-degree}\,(2)
to rewrite \(K_{\mathcal X}(\theta,u)\) as
\begin{equation}\label{eq:ch23-curve-polyhedron}
 K=\{v\in\mathbf R^I:\ v_i\leq u_i\ \forall i,\quad
     A_i(v)\geq0\ \forall i\},
\end{equation}
a closed convex polyhedron, since each \(A_i\) is affine.  Note that
\(\varphi_v\leq u\) as functions on \(X^{\An}\) is equivalent to
\(v_i\leq u_i\) for all \(i\): both functions are determined on
\(\mathcal X\), hence constant on the fibres of \(p_{\mathcal X}\) and
affine on each edge of \(\Delta_{\mathcal X}\), so the inequality on the
vertex set propagates to \(\Delta_{\mathcal X}\) and thence to \(X^{\An}\).

\emph{Step 1: \(K\ne\varnothing\).}  For a constant vector \(v\equiv t\) one
has \(A_i(t)=c_i\geq0\), because \(\theta\) is semipositive; and
\(t\leq\min_iu_i\) makes the obstacle constraint hold.  So
\(K\) contains all sufficiently negative constants.

\emph{Step 2: \(K\) is stable under componentwise maxima.}  Let
\(v,v'\in K\) and put \(w:=\max\{v,v'\}\) componentwise.  Fix \(i\) and
suppose without loss of generality \(w_i=v_i\).  Since \(w_j\geq v_j\) for
every \(j\) and \(w_{ij}\geq0\),
\[
 A_i(w)=c_i+\sum_{j\ne i}w_{ij}(w_j-w_i)
       \geq c_i+\sum_{j\ne i}w_{ij}(v_j-v_i)
       =A_i(v)\geq0 .
\]
Also \(w_i\leq u_i\).  Hence \(w\in K\).  (This is the graph-theoretic
shadow of \Cref{prop:ch20-model-max}; here it is a one-line computation
because the constraints are linear.)

\emph{Step 3: existence of a largest element.}  For each \(i\) set
\(P_i:=\sup\{v_i:v\in K\}\), which is finite because \(v_i\leq u_i\) and
\(K\ne\varnothing\).  Given \(\varepsilon>0\), choose for each \(i\) an
element \(v^{(i)}\in K\) with \(v^{(i)}_i>P_i-\varepsilon\), and let
\(w^{(\varepsilon)}:=\max_iv^{(i)}\), which lies in \(K\) by Step~2 and
satisfies \(P_i-\varepsilon<w^{(\varepsilon)}_i\leq P_i\) for every \(i\).
Letting \(\varepsilon\downarrow0\) and using that \(K\) is closed gives
\(P\in K\).  By construction \(P\geq v\) for every \(v\in K\), so \(P\) is
the largest element.  This proves (1).

\emph{Step 4: complementary slackness.}  Suppose, for contradiction, that
there is an index \(i\) with
\[
 P_i<u_i\qquad\text{and}\qquad A_i(P)>0 .
\]
Let \(\one_i\in\mathbf R^I\) be the \(i\)-th standard basis vector and, for
\(t>0\), put \(P^{(t)}:=P+t\one_i\).  We check that \(P^{(t)}\in K\) for
small \(t>0\):
\begin{itemize}
\item obstacle: \(P^{(t)}_i=P_i+t\leq u_i\) as soon as
\(t\leq u_i-P_i\), and \(P^{(t)}_j=P_j\leq u_j\) for \(j\ne i\);
\item the constraint at \(i\): with
\(s_i:=\sum_{j\ne i}w_{ij}\geq0\),
\[
 A_i(P^{(t)})=c_i+\sum_{j\ne i}w_{ij}(P_j-P_i-t)=A_i(P)-t\,s_i\geq0
\]
as soon as \(t\,s_i\leq A_i(P)\), which holds for small \(t>0\) because
\(A_i(P)>0\);
\item the constraints at \(j\ne i\): only the term with index \(i\) changes,
and it changes by \(+t\,w_{ji}\geq0\), so
\[
 A_j(P^{(t)})=A_j(P)+t\,w_{ji}\geq A_j(P)\geq0 .
\]
\end{itemize}
Thus \(P^{(t)}\in K\) and \(P^{(t)}>P\) in the \(i\)-th coordinate,
contradicting maximality of \(P\).  Hence for every \(i\) either
\(P_i=u_i\) or \(A_i(P)=0\), which is
\eqref{eq:ch23-complementary-slackness}.

\emph{Step 5: the integral form.}  By
\Cref{lem:ch23-curve-residual-degree}\,(3),
\[
 \int_{X^{\An}}(u-P)\,\MA_\theta(P)
 =\sum_{i\in I}A_i(P)\,(u_i-P_i)=0
\]
by \eqref{eq:ch23-complementary-slackness}.  Since \(u-P\geq0\) and
\(\MA_\theta(P)\geq0\), vanishing of the integral says exactly that the
measure gives no mass to \(\{u>P\}\), which in this atomic setting is the
set of vertices with \(P_i<u_i\).  This proves (2) and (3).
\end{proof}

\begin{corollary}
\label{cor:ch23-curve-harmonic-off-contact}
In the situation of \Cref{thm:ch23-model-orthogonality-curve}, let
\(N:=\{i\in I: P_i<u_i\}\) be the noncontact vertex set.  Then
\(A_i(P)=0\) for every \(i\in N\); equivalently, the restriction of
\(\MA_\theta(P)\) to the Borel set \(p_{\mathcal X}^{-1}(\{e_i:i\in N\})\)
vanishes, and on the corresponding part of the dual graph the vertical
divisor of \(P\) satisfies the homogeneous Laplace equation
\(\Delta_\Gamma P=-c\) exactly.
\end{corollary}

\begin{proof}
Immediate from \eqref{eq:ch23-complementary-slackness} and
\eqref{eq:ch23-curve-model-measure}; the last reformulation is
\eqref{eq:ch08-curve-cl-perturbed}, which identifies \(A_i(P)-c_i\) with
the graph Laplacian \eqref{eq:graph-laplacian} of \(P\) at \(e_i\).
\end{proof}

\begin{remark}[Why Step~4 does not survive in dimension \(\ge2\)]
\label{rem:ch23-model-slackness-fails}
Two features of dimension one were used, and both fail for \(n\geq2\).

First, the description of semipositivity.  In
\Cref{lem:ch23-curve-residual-degree}\,(2) the nef cone was cut out by the
\(\#I\) linear inequalities \(A_i(v)\geq0\), because in a one-dimensional
special fibre the only proper integral vertical curves are the components
themselves.  For \(n\geq2\) the special fibre contains a continuum of
vertical curves lying inside the components, and relative nefness is the
requirement that \emph{all} of them have nonnegative degree.  Raising one
coefficient \(v_i\) can then violate the condition on a curve
\(C\subset E_j\) with \(j\ne i\), even though it improves the degree on
\(E_j\) itself.

Second, the description of the measure.  The atom of the model measure at
\(e_i\) is \(b_i\bigl((\theta+\ddc\varphi)^n\cdot E_i\bigr)\), a
\emph{degree \(n\)} polynomial in the vertex values, not the affine function
\eqref{eq:ch23-curve-residual-degree}.  Positivity of that atom therefore
carries no information about the sign of a directional derivative of the
constraint set.

The correct higher-dimensional statement, and its genuine proof through
approximate Zariski decompositions and differentiability of non-Archimedean
volumes, is \Cref{thm:orthogonality-provenance}; the warning in
\Cref{rem:model-slackness-heuristic} is precisely the observation that the
argument above is one dimensional.
\end{remark}

\subsection{The polyhedral convex case}
\label{sec:ch23-convex-orthogonality}

The second elementary case is the invariant toric one, where the obstacle
problem becomes the classical convex-envelope problem and the
Monge--Amp\`ere measure becomes the Alexandrov measure of
\Cref{sec:appD-real-ma}.  Here orthogonality holds in every dimension, and
the proof is a Legendre-duality argument.  We state it as a theorem of
convex analysis; the translation to invariant metrics on a polarized toric
variety is \Cref{ex:ch23-toric-orthogonality}.

\begin{theorem}
\label{thm:ch23-convex-envelope-orthogonality}
Let \(P\subset\mathbf R^n\) be a compact convex body and let
\(f:P\to\mathbf R\) be continuous.  Let
\[
 u:=\sup\{\ell:\ \ell\text{ affine},\ \ell\leq f\text{ on }P\}
\]
be the convex envelope of \(f\), and put
\(C:=\{x\in P^\circ:u(x)=f(x)\}\).  Then \(C\) is closed in \(P^\circ\) and
\begin{equation}\label{eq:ch23-convex-orthogonality}
 \MA_{\mathbf R}(u)\bigl(P^\circ\setminus C\bigr)=0 .
\end{equation}
\end{theorem}

\begin{proof}
The constant \(\min_Pf\) is an affine minorant, so \(u\) is finite.  It is
convex as a supremum of affine functions and is continuous on \(P^\circ\);
hence \(C\) is closed there.

\emph{Step 1: \(u\) is affine along a segment through each noncontact
point.}  Let \(x\in P^\circ\setminus C\).  The barycentric formula for the
convex envelope on a compact convex body gives points
\(y_0,\ldots,y_q\in P\), with \(q\leq n\), and positive weights
\(\lambda_k\) such that
\(\sum_k\lambda_k=1\), \(\sum_k\lambda_ky_k=x\), and
\begin{equation}\label{eq:ch23-caratheodory}
                         u(x)=\sum_k\lambda_kf(y_k).
\end{equation}
Indeed, minimize \(\sum_k\lambda_kf(y_k)\) over probability measures on
\(P\) with barycentre \(x\).  Compactness gives a minimizer, and the
finite-dimensional moment constraints admit one supported at at most
\(n+1\) points.  Separating its epigraph shows that the minimum equals the
supremum of the affine minorants, hence equals \(u(x)\).  Choose a minimizer
with minimal support; its support is affinely independent, and \(x\) lies in
the relative interior of its convex hull.

Convexity and \(u\leq f\) give
\[
 u(x)\leq\sum_k\lambda_ku(y_k)
       \leq\sum_k\lambda_kf(y_k)=u(x).
\]
Thus both inequalities are equalities, and positivity of the weights gives
\(u(y_k)=f(y_k)\) for every \(k\).  Since \(x\) is not a contact point, the
support contains at least two points.  Let \(p\in\partial u(x)\) be
arbitrary; this set is nonempty because \(x\in P^\circ\).  Equality above forces
\(u(y_k)=u(x)+\langle p,y_k-x\rangle\) for every \(k\).  Hence \(u\) agrees
with this affine support function on
\(S:=\operatorname{conv}\{y_0,\ldots,y_q\}\), which contains a nontrivial
segment through \(x\).

\emph{Step 2: subgradients at noncontact points are singular for the
conjugate.}  Extend \(u\) by \(+\infty\) off \(P\), and let \(u^*\) be its
Legendre conjugate.  Since \(P\) is compact and
\(\min_Pf\leq u\leq\max_Pf\), the function \(u^*\) is finite on
\(\mathbf R^n\).  For \(x\in P^\circ\setminus C\), Step~1 shows that every
point of a nontrivial segment through \(x\) realizes equality in the
Fenchel inequality for \((u,u^*)\) and the same \(p\).  Thus this segment is
contained in \(\partial u^*(p)\) by
\eqref{eq:appD-subgradient}; in particular \(u^*\) is not differentiable at
\(p\).

\emph{Step 3: conclusion.}  The set of points at which the finite convex
function \(u^*\) fails to be differentiable has Lebesgue measure zero: this
is the standard differentiability theorem for convex functions used in the
proof of \Cref{thm:appD-alexandrov-measure}.  By Step~2,
\(\partial u(P^\circ\setminus C)\) is contained in that set, so
\[
 \MA_{\mathbf R}(u)(P^\circ\setminus C)
 =\operatorname{Leb}\bigl(\partial u(P^\circ\setminus C)\bigr)=0
\]
by the definition \eqref{eq:appD-alexandrov-measure} of the Alexandrov
measure.
\end{proof}

\begin{remark}[What the convex case does and does not prove]
\label{rem:ch23-convex-case-scope}
\Cref{thm:ch23-convex-envelope-orthogonality} is a theorem about convex
functions on a compact convex body, and it transfers to invariant metrics on a
polarized toric variety because there the Monge--Amp\`ere operator
\emph{is} the real Monge--Amp\`ere operator, up to the factor \(n!\)
recorded in \Cref{app:normalizations}, and the psh envelope \emph{is} the
convex envelope; see \Cref{ex:ch23-toric-orthogonality}.  It is not a proof
of \eqref{eq:ch23-orthogonality-preview} for a general model, for the reason
given in \Cref{rem:ch23-model-slackness-fails}: outside the invariant
setting there is no linear or convex parameterization of the psh cone, and
the passage from ``\(u\) is affine along a segment'' to ``the measure does
not charge the point'' has no counterpart.
\end{remark}

\subsection{Summary of the model-level results}
\label{sec:ch23-orthogonality-summary}

\begin{center}
\small
\begin{tabular}{p{.24\textwidth}|p{.30\textwidth}|p{.36\textwidth}}
\textbf{Setting}&\textbf{Statement proved here}&\textbf{Mechanism}\\ \hline
curve model, model obstacle
 &\Cref{thm:ch23-model-orthogonality-curve}
 &linear constraints; raising one coordinate stays in the cone\\
polytope, continuous obstacle
 &\Cref{thm:ch23-convex-envelope-orthogonality}
 &Legendre duality; nondifferentiability set of \(u^*\) is null\\
general model, \(n\geq2\)
 &not proved here
 &requires energy, volumes and asymptotic orthogonality;
  \Cref{thm:orthogonality-provenance}\\
\end{tabular}
\end{center}

\section{Worked examples}
\label{sec:ch23-worked}

\begin{example}[A globally generated reference model]
\label{ex:ch23-globally-generated}
Take \(X=\mathbf P^n_K\), its standard model
\(\mathcal X=\mathbf P^n_R\), and
\(\mathcal L=\mathcal O_{\mathbf P^n_R}(1)\).  Every
\(m\mathcal L\) is globally generated, so
\[
                 \mathfrak a_m=\mathcal O_{\mathcal X},
                 \qquad \varphi_m=0.
\]
The zero function is \(\theta\)-psh and lies below the zero obstacle, while
every competitor does too.  Hence \(P_\theta(0)=0\), and
\eqref{eq:ch23-base-ideal-convergence} is an equality at every finite level;
in \Cref{cor:ch23-envelope-rate} one may take \(C=0\).  The
multiplier-ideal correction becomes necessary only when the chosen
model extension has vertical base locus.  For an example where it does, keep
\(X=\mathbf P^1_K\) and replace \(\mathcal X\) by the blowup
\(\mathcal X_1\) of \Cref{ex:ch23-pole-model-check} with
\(\mathcal L_1:=\pi^*\cO(1)-E_1\), whose generic restriction is
\(\cO(1)\).  Then \(H^0(\mathcal X_1,m\mathcal L_1)\) consists of those
sections of \(m\cO(1)\) vanishing to order \(m\) along \(E_1\), and the
resulting base ideal is nontrivial; the corresponding \(\varphi_m\) is
\(-1\) at \(e_1\) for every \(m\), and the envelope \(P_\theta(0)\) is the
function \(\psi_1=\max\{\psi,-1\}\) of
\eqref{eq:ch23-truncated-pole}.
\end{example}

\begin{example}[A convex obstacle on an interval]
\label{ex:ch23-convex-interval}
On a one-dimensional skeleton, after subtracting a fixed reference
potential, the local obstacle problem becomes the largest convex minorant.
For
\begin{equation}\label{eq:ch23-concave-obstacle}
                         f(t)=t(1-t),
                         \qquad 0\leq t\leq1,
\end{equation}
any convex \(u\leq f\) satisfies
\[
                  u(t)\leq(1-t)u(0)+tu(1)\leq0 .
\]
Since the zero function is admissible, the convex envelope is identically
zero.  It touches \(f\) only at the endpoints and is affine on the
noncontact interval.  In the metric-graph realization, the variable
Laplacian consequently vanishes off the contact set, in accordance with
\Cref{cor:ch23-curve-harmonic-off-contact} and
\Cref{thm:ch23-convex-envelope-orthogonality}: the contact set is
\(\{0,1\}\), the measure is \(\delta_0+\delta_1\) up to the reference term,
and \(\int(f-u)\,\MA_{\mathbf R}(u)=f(0)\cdot1+f(1)\cdot1=0\).  The fixed
reference curvature must still be restored when translating this local
calculation to a global polarized curve; that is done in
\Cref{ex:ch23-curve-envelope-numbers}.
\end{example}

\begin{example}[An obstacle envelope on a curve skeleton, in numbers]
\label{ex:ch23-curve-envelope-numbers}
This is the promised complete computation.  Let \(K\) be discretely valued
with \(v(\pi)=1\), let \(X=\mathbf P^1_K\), and let \(\mathcal X\) be the
model obtained from \(\mathbf P^1_R\) by four successive blowups, each at a
smooth point of the special fibre lying on the component created by the
previous blowup and away from the previous ones.  Then \(\mathcal X\) is
regular with reduced special fibre a chain
\[
 \mathcal X_s=E_0+E_1+E_2+E_3+E_4,
 \qquad b_i=1,
\]
with \((E_i\cdot E_{i+1})=1\), all other products zero, and
\[
 (E_0^2)=(E_4^2)=-1,\qquad (E_1^2)=(E_2^2)=(E_3^2)=-2 ,
\]
the row sums \((\mathcal X_s\cdot E_i)\) being zero as they must.  The dual
graph \(\Gamma\) is the path \(e_0-e_1-e_2-e_3-e_4\), and in the notation
\eqref{eq:ch23-curve-special-fibre} all edge weights are
\(w_{i,i+1}=1\).

Take \(L=\cO_{\mathbf P^1}(2)\), of degree \(2\), and let \(\mathcal L\) be
the extension of \(L\) to \(\mathcal X\) with residual degrees
\[
 (c_0,c_1,c_2,c_3,c_4)=(1,0,0,0,1),
 \qquad \sum_ic_i=2=\deg_XL .
\]
Such an \(\mathcal L\) exists.  Start from any extension and twist by
vertical divisors: twisting by \(E_0\) changes the degree vector by
\(\bigl((E_0\cdot E_i)\bigr)_i=(-1,1,0,0,0)\), twisting by \(E_0+E_1\)
changes it by \((0,-1,1,0,0)\), and in general twisting by
\(E_0+\cdots+E_j\) moves one unit of degree from \(e_j\) to \(e_{j+1}\).
These moves generate all integer vectors of total sum zero, and the total
degree \(\sum_ic_i=\deg_XL\) is fixed by
\eqref{eq:ch23-curve-degree-sum}; so every prescribed vector with the
correct total degree is realized.  Put
\(\theta=c_1(\mathcal L,\|\cdot\|_{\mathcal L})\); it is semipositive
because all \(c_i\geq0\).

For the obstacle take the model function \(u\) determined on \(\mathcal X\)
with vertex values
\[
 (u_0,u_1,u_2,u_3,u_4)=(0,-\tfrac32,-3,-\tfrac32,0),
\]
that is, the piecewise affine ``V'' with vertex value \(-3\) at \(e_2\),
affine on \([e_0,e_2]\) and on \([e_2,e_4]\).

\emph{The constraints.}  By
\Cref{lem:ch23-curve-residual-degree} the competitor set is
\[
 K=\Bigl\{v\in\mathbf R^5:\ v_i\leq u_i,\quad
   A_0=1+(v_1-v_0)\geq0,\ \
   A_4=1+(v_3-v_4)\geq0,
\]
\[
 \phantom{K=\Bigl\{}
   A_i=(v_{i-1}-v_i)+(v_{i+1}-v_i)\geq0\ \ (i=1,2,3)\Bigr\},
\]
the last three conditions being convexity of \(v\) along the path.

\emph{The envelope.}  We compute \(P=P^{\mathcal X}_\theta(u)\) coordinate
by coordinate, using that \(P_i=\max\{v_i:v\in K\}\) by Step~3 of the proof
of \Cref{thm:ch23-model-orthogonality-curve}.  First \(P_2\leq u_2=-3\), and
\(P_2=-3\) is attained, as the explicit vector below shows.  Next, for
\(v\in K\), convexity at \(e_1\) gives \(v_0+v_2\geq2v_1\) and \(A_0\geq0\)
gives \(v_0\leq v_1+1\); combining,
\[
 2v_1\leq v_0+v_2\leq (v_1+1)+(-3),
 \qquad\text{so}\qquad v_1\leq-2 ,
\]
and then \(v_0\leq v_1+1\leq-1\).  By symmetry \(v_3\leq-2\) and
\(v_4\leq-1\).  All four bounds are attained simultaneously by
\[
 \boxed{\ P=(-1,-2,-3,-2,-1)\ }
\]
which indeed lies in \(K\): \(P_i\leq u_i\) holds coordinatewise
(\(-1\leq0\), \(-2\leq-\tfrac32\), \(-3\leq-3\)), and
\[
 A_0(P)=1+(-2+1)=0,\quad
 A_1(P)=(-1+2)+(-3+2)=0,
\]
\[
 A_2(P)=(-2+3)+(-2+3)=2,\quad
 A_3(P)=0,\quad
 A_4(P)=0 .
\]

\emph{The measure, the mass and orthogonality.}  By
\eqref{eq:ch23-curve-model-measure},
\[
 \MA_\theta(P)=2\,\delta_{e_2},
 \qquad
 \MA_\theta(P)(X^{\An})=2=\deg_XL ,
\]
so the total mass is correct.  The contact set is \(\{e_2\}\), since
\(P_2=u_2=-3\) while \(u_i-P_i=1,\ \tfrac12,\ 0,\ \tfrac12,\ 1\) for
\(i=0,\ldots,4\).  Therefore
\[
 \int_{X^{\An}}(u-P)\,\MA_\theta(P)
 =1\cdot0+\tfrac12\cdot0+0\cdot2+\tfrac12\cdot0+1\cdot0=0 ,
\]
which is \eqref{eq:ch23-model-orthogonality} verified in numbers.  Observe
the two features predicted by
\Cref{cor:ch23-curve-harmonic-off-contact}: the four noncontact vertices
carry no mass, and \(P\) is affine of slope \(-1\) on \([e_0,e_2]\) and of
slope \(+1\) on \([e_2,e_4]\), that is, harmonic off the contact set, with
the two boundary slopes exactly saturating the constraints
\(A_0=A_4=0\) coming from the residual degrees \(c_0=c_4=1\).

\emph{\(P\) is the global envelope.}  Finally we check that no refinement of
\(\mathcal X\) improves \(P\), so that \(P=P_\theta(u)\) on
\(\Delta_{\mathcal X}\).  Let \(\varphi\in\PSH(X,\theta)\) with
\(\varphi\leq u\).  Its restriction to \(\Delta_{\mathcal X}\) is convex on
each edge by \Cref{thm:bfj-skeletal-characterization}\,(1).  We claim its
outgoing slope at \(e_0\) along the path, measured in the unit-length
parameterization of the edges, is \(\geq-c_0=-1\).  Granting the claim,
convexity along the path makes all subsequent slopes \(\geq-1\), so
\[
 \varphi(e_0)\leq\varphi(e_2)+2\leq u_2+2=-1,
\]
and then \(\varphi(e_1)\leq\varphi(e_2)+1=-2\); together with the symmetric
estimates this gives \(\varphi\leq P\) on the vertex set, hence on
\(\Delta_{\mathcal X}\).  The claim itself is the slope criterion
\eqref{eq:curve-slope-criterion}, which for a \(\theta\)-psh function on a
regular semistable model reads: at every point of the skeleton, the sum of
outgoing slopes plus the residual degree of \(\theta\) on the corresponding
component is nonnegative.  At \(e_0\) the only outgoing direction inside
\(\Delta_{\mathcal X}\) is towards \(e_1\), and the directions leaving the
skeleton contribute nonpositive slopes by the retraction inequality, so
discarding them preserves the inequality; see the proof of
\Cref{thm:curves-arbitrary-field-continuity}, where this reduction is
carried out in detail.
\end{example}

\begin{remark}[Model dependence of the model-level envelope]
\label{rem:ch23-model-dependence}
The chain \eqref{eq:ch23-model-level-below} is a chain of inequalities, and
in general there is no reason for it to stabilize at a finite level:
\Cref{def:ch23-fixed-model-envelope} constrains the competitors to be affine
on the faces of \(\Delta_{\mathcal X}\), and refining the model enlarges the
competitor set.  The finite-dimensional shadow of
\Cref{ex:metrics-max-skeleton} is that a refinement can create a new vertex
at which a competitor may be raised.

In dimension one, however, the chain does stabilize, and it stabilizes as
soon as the model sees the data.  Suppose \(\mathcal X\) is a strictly
semistable model, available after a finite extension of \(K\), determining
both \(\theta\) and the model obstacle \(u\).  The chord argument in the
proof of \Cref{thm:curves-arbitrary-field-continuity}---replace a competitor
by the piecewise affine function agreeing with it at the vertices, which can
only increase the outgoing slopes and hence preserves the slope criterion
\eqref{eq:curve-slope-criterion}---shows that the supremum defining
\(P_\theta(u)\) is unchanged if one restricts to competitors affine on the
edges of \(\Delta_{\mathcal X}\).  Hence
\[
 P^{\mathcal X}_\theta(u)=P_\theta(u)
 \qquad(n=1,\ \mathcal X\text{ strictly semistable and determining }
 \theta,u),
\]
which is the equality verified by hand at the end of
\Cref{ex:ch23-curve-envelope-numbers}.  This is one more item on the list of
\Cref{rem:ch23-curve-specialness}: it uses that the psh condition is linear
on each skeleton and that retraction preserves it, and it has no known
higher-dimensional analogue.
\end{remark}

\begin{example}[Toric convex envelopes and orthogonality in numbers]
\label{ex:ch23-toric-orthogonality}
For an invariant continuous metric on a polarized toric variety, subtracting
the reference support function converts the psh envelope into the convex
envelope with prescribed recession behavior.  Uniform convergence in
\eqref{eq:ch23-envelope-continuity} becomes uniform convergence of convex
piecewise affine minorants.  The Alexandrov measure is supported on the
contact set by \Cref{thm:ch23-convex-envelope-orthogonality}, so
orthogonality can be verified directly in this invariant setting.  This
direct convex proof is an instance of
\eqref{eq:ch23-orthogonality-preview}; it is not a proof for arbitrary
higher-dimensional models, by \Cref{rem:ch23-convex-case-scope}.

Here is the computation, first in one variable and then in two.  On
\(P_1=[-2,2]\) take the continuous, nonconvex obstacle
\[
 \psi(t):=\min\{|t|,1\} .
\]
Its convex envelope is
\[
 u_1(t)=\tfrac12|t| .
\]
Indeed \(u_1\) is convex and \(u_1\leq\psi\), since
\(\tfrac12|t|\leq|t|\) for \(|t|\leq1\) and \(\tfrac12|t|\leq1\) for
\(|t|\leq2\); conversely if \(u\leq\psi\) is convex then \(u(0)\leq0\) and
\(u(\pm2)\leq1\), so for \(0\leq t\leq2\)
\[
 u(t)\leq\Bigl(1-\tfrac t2\Bigr)u(0)+\tfrac t2\,u(2)\leq\tfrac t2 ,
\]
and symmetrically for \(t\leq0\).  The contact set is
\(C_1=\{-2,0,2\}\), because \(\tfrac12|t|=\psi(t)\) forces \(t=0\) or
\(|t|=2\).  The Alexandrov measure of \(u_1\) is the jump of its derivative
at the origin,
\[
 \MA_{\mathbf R}(u_1)=\bigl(\tfrac12-(-\tfrac12)\bigr)\delta_0=\delta_0 ,
\]
carried by \(C_1\), and
\(\int(\psi-u_1)\,\MA_{\mathbf R}(u_1)=(\psi(0)-u_1(0))\cdot1=0\).

Now let \(n=2\), \(P=P_1^2=[-2,2]^2\), and take the separable obstacle
\[
 f(x_1,x_2):=\psi(x_1)+\psi(x_2).
\]
The convex envelope of a separable sum is the sum of the convex envelopes:
the sum \(u_1(x_1)+u_1(x_2)\) is convex and dominated by \(f\), and any
convex \(u\leq f\) is dominated by it because one may take convex envelopes
in each variable separately.  Hence
\[
 u(x_1,x_2)=\tfrac12\bigl(|x_1|+|x_2|\bigr),
 \qquad
 C=\{-2,0,2\}\times\{-2,0,2\},
\]
the contact set being a product because equality in a sum of two
inequalities forces equality in each.  The subdifferential of \(u\) at the
origin is the square \([-\tfrac12,\tfrac12]^2\), and at every other point it
is a face of that square; so by \eqref{eq:appD-alexandrov-measure} and
\Cref{ex:appD-polyhedral-atom}
\[
 \MA_{\mathbf R}(u)=\operatorname{Leb}\bigl([-\tfrac12,\tfrac12]^2\bigr)
   \,\delta_{(0,0)}=1\cdot\delta_{(0,0)},
\]
of total mass \(1=\operatorname{Leb}(\partial u(P^\circ))\), in accordance
with \eqref{eq:appD-real-ma-mass}.  The measure is carried by the single
point \((0,0)\in C\), so
\[
 \int_{P^\circ}(f-u)\,\MA_{\mathbf R}(u)
 =\bigl(f(0,0)-u(0,0)\bigr)\cdot1=0 ,
\]
which is \eqref{eq:ch23-convex-orthogonality} in numbers.  Note that the
noncontact region here has full Lebesgue measure and yet carries no
Monge--Amp\`ere mass at all: the envelope is affine along a segment through
each of its points, exactly as in Step~1 of the proof of
\Cref{thm:ch23-convex-envelope-orthogonality}.  Translating through the
toric dictionary, \(u\) is the support-function description of an invariant
semipositive metric whose Chambert--Loir measure is \(n!=2\) times the atom
above, by the \(n!\)-test recorded in \Cref{app:normalizations}.
\end{example}

\begin{example}[Stability under perturbing the obstacle]
\label{ex:ch23-obstacle-stability}
If \(f_j\to f\) uniformly, then
\eqref{eq:ch23-envelope-contraction} gives
\[
                   \|P_\theta(f_j)-P_\theta(f)\|_\infty
                   \leq\|f_j-f\|_\infty\longrightarrow0.
\]
Thus all analytic difficulty lies in proving continuity for a dense class of
obstacles, not in passing from that class to every continuous obstacle.  The
base-ideal theorem treats model obstacles, and the contraction estimate
performs the passage.  The same remark identifies the exact reason that
\Cref{lem:ch23-reduction-to-base-ideals} is a reduction and not a
restriction: model functions are dense by
\Cref{lem:ch23-model-density}, and the envelope operator is a sup-norm
contraction, so a statement proved for a dense set of obstacles is proved for
all of them.  Note that this argument transports \emph{continuity} but not
\emph{a rate}: the modulus of continuity of \(P_\theta(f)\) is not controlled
by \eqref{eq:ch23-envelope-contraction}.
\end{example}

\begin{example}[An envelope outside the scope of every theorem proved or
cited here]
\label{ex:ch23-outside-scope}
Let \(K=\mathbf Q_p\), which is complete and discretely valued with residue
field \(\mathbf F_p\), so of \emph{mixed} characteristic \((0,p)\).  Let
\(X\) be a smooth projective threefold over \(K\) with an ample line bundle
\(L\), for instance a smooth hypersurface of degree \(5\) in
\(\mathbf P^4_{\mathbf Q_p}\) with \(L=\cO(1)|_X\), and let
\(\theta=c_1(L,\|\cdot\|_0)\) for a model metric.  Let \(f\in C^0(X^{\An})\)
be arbitrary.

None of the theorems available to us applies to \(P_\theta(f)\).
\Cref{thm:bfj-envelope-regularity} requires residue characteristic zero,
which fails, and its proof fails twice over: SNC models need not be cofinal
(\Cref{rem:cofinality-scope}), so even
\Cref{def:ch21-bfj-psh} would have to be reconsidered, and asymptotic
multiplier ideals with subadditivity and fixed-twist global generation are
not available.  The test-ideal theorem
\Cref{thm:ch23-gjkm-positive-characteristic} requires \emph{equal}
characteristic \(p\) and geometric origin from a family over a perfect
field, both of which fail for \(\mathbf Q_p\).  The curve theorem
\Cref{thm:curves-arbitrary-field-continuity} requires \(\dim X=1\).  And the
removal of algebraizability in
\Cref{thm:bggjkm-remove-algebraizability} presupposes envelope continuity
rather than proving it.

The honest statement is therefore: for a smooth projective variety of
dimension at least three over a mixed-characteristic complete discretely
valued field, continuity of \(P_\theta(f)\) for a general continuous \(f\)
does not follow from anything proved or cited in this chapter.  Partial
results in mixed characteristic, and the precise form of the
Monge--Amp\`ere theory that is currently available there, are recorded in
\Cref{thm:fgk-mixed-characteristic-cy} with their own hypotheses and date.
We do not assert that continuity fails; we assert that it is not known by
these methods, and the distinction is exactly the one that
\Cref{sec:ch23-later} is designed to keep visible.
\end{example}

\section{Later envelope theorems and their exact extra hypotheses}
\label{sec:ch23-later}

The multiplier-ideal proof above belongs to the original complete discretely
valued, residue-characteristic-zero setting.  Four later developments extend
it, and each of them buys generality in one direction by paying in another.
We state them separately, with their own hypotheses, sources and dates,
rather than folding them into
\Cref{thm:bfj-envelope-regularity}.  The corresponding
Monge--Amp\`ere consequences, which need more than envelope continuity, are
in \Cref{chap:bfj-continuity}; here we record only what concerns envelopes.

\subsection{Curves over an arbitrary complete field}
\label{sec:ch23-later-curves}

The first extension removes discreteness and residue-characteristic
restrictions for ample line bundles on curves.

\begin{theorem}
\label{thm:ch23-gjkm-curves}
Let \(K\) be complete with respect to a nontrivial non-Archimedean
absolute value, let \(X\) be a smooth projective geometrically connected
curve over \(K\), and let \(L\) be ample.  For every continuous metric
on \(L^{\An}\), its envelope over semipositive rational model metrics
is a continuous semipositive metric.
\end{theorem}

\begin{proof}
This is \cite[Theorem~1.1]{GJKM19}.  In potential notation it is the
rational-model envelope, not by definition the raw all-psh envelope.
Ampleness is retained; no existence of a rational model competitor is
inferred here from generic nefness alone.
\end{proof}

\begin{remark}[Why dimension one is special, and what fails above it]
\label{rem:ch23-curve-specialness}
Two things are true on a curve and false in general.  First, the
Monge--Amp\`ere operator is the graph Laplacian, so the psh condition is a
system of \emph{linear} inequalities on each skeleton; this is the same
feature exploited in \Cref{thm:ch23-model-orthogonality-curve}.  Second, the
retraction of a semipositive model metric to a skeleton is again
semipositive: \cite[Proposition~1.2 and Propositions~3.5, 3.8]{GJKM19}.  The
second statement is exactly what fails in higher dimension.  Burgos Gil and
Sombra, in the appendix to \cite{GJKM19}, exhibit two toric models of a
toric surface with the same skeleton but different retractions, and deduce
that the retracted metric \(\|\cdot\|_\Delta\) of a semipositive model
metric need not be semipositive.  Consequently the combinatorial reduction
``work on one skeleton'' is unavailable for \(n\geq2\), which is why the
proof of \Cref{thm:bfj-envelope-regularity} has to leave the skeleta and use
ideals on models.
\end{remark}

\subsection{Equal positive characteristic: test ideals}
\label{sec:ch23-later-test-ideals}

The second extension keeps arbitrary dimension and replaces multiplier
ideals by test ideals.  Three hypotheses appear that have no analogue in
\Cref{thm:bfj-envelope-regularity}, and all three are genuine.

\begin{definition}[Geometric origin; {\cite[Definition~8.1]{GJKM19}}]
\label{def:ch23-geometric-origin}
Let \(K\) be complete discretely valued of characteristic \(p>0\).  One says
that \(X\) is \emph{of geometric origin from a \(d\)-dimensional family over
a field \(k\)} if there are a \(d\)-dimensional normal variety \(B\) over
\(k\), a codimension-one point \(b\in B^{(1)}\), a subfield
\(K'\subset K\) inducing \(K\) as the completion of the discrete valuation
ring \(\cO_{B,b}\), and a \(K'\)-variety \(Y\) with
\(Y\otimes_{K'}K\simeq X\).  If in addition a line bundle, respectively a
closed \((1,1)\)-form, on \(X\) is induced by base change from that family,
one says that \((X,L)\), respectively \((X,\theta)\), is of geometric
origin.
\end{definition}

\begin{theorem}
\label{thm:ch23-gjkm-positive-characteristic}
Let \(K\) be complete discretely valued of characteristic \(p>0\), let \(X\)
be a smooth projective \(K\)-variety of dimension \(n\) which is of
geometric origin from a \(d\)-dimensional family over a perfect field
\(k\), and assume that
\begin{center}
\textup{(R)}\quad resolution of singularities holds over \(k\)
in dimension \(d+n\).
\end{center}
Let \(\theta\) be a closed \((1,1)\)-form with ample de Rham class and let
\(u\in C^0(X^{\An})\).  Then \(P_\theta(u)\) is a uniform limit of
\(\theta\)-psh model functions, and in particular is continuous.
\end{theorem}

\begin{proof}[Source and the shape of the proof]
This is \cite[Theorem~8.2]{GJKM19}.  The architecture is that of
\Cref{thm:bfj-envelope-regularity}, with
\Cref{thm:ch23-uniform-generation-input} replaced throughout by the
corresponding statements for asymptotic test ideals
\(\tau(\lambda\|D\|)\) of \Cref{sec:appC-test-ideals}: containment,
subadditivity, and a uniform generation theorem
\cite[Theorem~1.3 and Theorem~4.6]{GJKM19}, which asserts that for
\(X\) projective over a finitely generated \(k\)-algebra \(R\) and smooth of
dimension \(n\) over \(k\), for \(H\) ample and basepoint free, \(D\) with
\(h^0(mD)\neq0\) for some \(m>0\), and \(E\) with \(D-\lambda E\) nef for
some \(\lambda\in\mathbf Q_{\geq0}\), the sheaf
\(\cO_X(K_{X/k}+E+dH)\otimes\tau(\lambda\|D\|)\) is globally generated for
all \(d\geq n+1\).  Mustaţă proved this for \(X\) projective over \(k\);
the version over a base is obtained in \cite[Theorem~4.6]{GJKM19} by
replacing Fujita vanishing with Keeler's generalization.  Assumption
\textup{(R)} is used to produce the regular models on which the test-ideal
machinery operates, and a semi-factoriality result of P\'epin enters at
\cite[Lemma~7.5]{GJKM19}.
\end{proof}

\begin{corollary}
\label{cor:ch23-gjkm-surfaces}
Let \(X\) be a smooth projective \emph{surface} over a complete discretely
valued field \(K\) of characteristic \(p>0\), of geometric origin from a
one-dimensional family over a perfect field \(k\).  Then the conclusion of
\Cref{thm:ch23-gjkm-positive-characteristic} holds unconditionally.
\end{corollary}

\begin{proof}[Source]
Here \(d+n=3\), and resolution of singularities in dimension three over a
perfect field is the theorem of Cossart--Piltant quoted as
\cite[Theorem~6.3]{GJKM19}.  Nothing else changes.
\end{proof}

\begin{remark}[Reading the three extra hypotheses]
\label{rem:ch23-gjkm-hypotheses}
\emph{Equal characteristic.}  The field \(K\) must have characteristic
\(p\), not merely residue characteristic \(p\); test ideals are defined by
Frobenius on \(K\)-schemes, and there is no Frobenius in mixed
characteristic.  This is why \Cref{ex:ch23-outside-scope} is outside the
scope of \Cref{thm:ch23-gjkm-positive-characteristic}.

\emph{Geometric origin.}  \Cref{def:ch23-geometric-origin} is what allows
the whole situation to be spread out over a variety \(B\) over a perfect
field, so that test ideals may be formed on a scheme of finite type over
\(k\) rather than over \(K^\circ\).  It is a restriction on \(X\), not a
normalization: not every smooth projective variety over
\(\mathbf F_p(\!(t)\!)\) is known to be of geometric origin from a family
over \(\mathbf F_p\).

\emph{Resolution.}  Assumption \textup{(R)} is used in dimension \(d+n\),
not \(n\); the family contributes its own dimension.  For \(n=2\),
\(d=1\) it is a theorem, whence
\Cref{cor:ch23-gjkm-surfaces}.  For \(n\geq3\) it is open in general.

Finally, in the setting of
\Cref{thm:ch23-gjkm-positive-characteristic} the two envelope definitions of
\Cref{sec:ch23-three-operators} are again identified, but only \emph{after}
the continuity theorem: \cite[Section~9]{GJKM19} deduces the agreement of
\Cref{def:ch23-envelope} and \Cref{def:ch23-model-envelope} from continuity,
exactly as \Cref{cor:ch23-envelope-by-model-functions} does here.  The
embedded-resolution assumption is then added, on top of \textup{(R)}, to
make projective SNC models cofinal and so to reconstruct the psh theory of
\Cref{chap:bfj-psh}; see \Cref{thm:gjkm-positive-characteristic-cy}.
\end{remark}

\subsection{Alterations, and what they do and do not replace}
\label{sec:ch23-later-alterations}

Both of the previous theorems need models with good singularities.  In
residue characteristic zero these are supplied by resolution, and
\Cref{prop:snc-cofinal} makes SNC models cofinal.  Outside that regime the
only general tool is de Jong's alteration theorem
\Cref{thm:appB-de-jong}, and it is important to see precisely what it does
not give.

An alteration is a proper, surjective, generically finite morphism
\(X'\to X\), not a modification; the induced map on skeleta is not an
inclusion, and there is no retraction from \(X^{\An}\) to a skeleton
produced upstairs.  Consequently one obtains, after a finite extension of
\(K\), semistable models of \(X'\) but not a cofinal system of SNC models of
\(X\); \Cref{rem:cofinality-scope} states the consequences.  For envelopes
this means:
\begin{enumerate}
\item for the rational-model envelope, pullback along a finite normal
extension \(K'/K\) commutes with the envelope by
\cite[Proposition~2.11]{GJKM19}.  One may enlarge a finite extension to
a finite normal extension; the cited proposition includes the purely
inseparable case.  This is not a statement about arbitrary alterations or,
without a separate comparison, the raw all-psh envelope;
\item pullback along a genuine alteration \(X'\to X\) does \emph{not}
obviously reflect continuity of envelopes, because the psh class on \(X\) is
not the set of functions whose pullback is psh unless a descent theorem is
available;
\item accordingly, no theorem in this book deduces
\Cref{thm:bfj-envelope-regularity} from de Jong's theorem.  Alterations are
used here for semistable reduction of curves, in
\Cref{thm:curves-arbitrary-field-continuity}, and for the model-theoretic
inputs of \Cref{chap:models}, not as a substitute for resolution in the
envelope proof.
\end{enumerate}
Descent of the psh property along finite and alteration maps is a serious
matter, and it is exactly the technique developed, in the trivially valued
setting, in \Cref{chap:bj-alterations}; the reader should not import it here
without checking that its hypotheses are met.

\subsection{Removing algebraizability, and mixed characteristic}
\label{sec:ch23-later-algebraizability}

Two further developments belong to the record, although neither is a
statement about continuity of envelopes.

Burgos Gil--Gubler--Jell--K\"unnemann--Martin proved differentiability of
non-Archimedean volumes and deduced the orthogonality property from
continuity of the envelope, thereby removing the algebraizability
hypothesis under which orthogonality had originally been proved in
\cite{BFJ15}; the precise statement, with its hypotheses, is
\Cref{thm:bggjkm-remove-algebraizability}, and the provenance ledger is
\Cref{thm:orthogonality-provenance}.  This does not weaken the hypotheses of
\Cref{thm:bfj-envelope-regularity}; it changes what may be deduced
\emph{from} that theorem.

In mixed characteristic the currently available Monge--Amp\`ere theory, and
the hypotheses under which it is available, are recorded in
\Cref{thm:fgk-mixed-characteristic-cy}.  \Cref{ex:ch23-outside-scope}
records an explicit instance in which continuity of the envelope is not
covered by any theorem proved or cited here.

Finally, the trivially valued unibranch theorem of
\Cref{chap:bj-envelope-statement} is a different result about a different
space, with a different proof: its hypotheses are a trivially valued
algebraically closed ground field, an integral projective unibranch variety
of arbitrary characteristic, and an arbitrary real numerical class, and its
conclusion concerns regularized suprema rather than continuity.  It neither
implies nor is implied by \Cref{thm:bfj-envelope-regularity}.

\subsection{Hypothesis atlas}
\label{sec:ch23-atlas}

\begin{center}
\small
\begin{tabular}{p{.19\textwidth}|p{.14\textwidth}|p{.13\textwidth}|p{.15\textwidth}|p{.24\textwidth}}
\textbf{Statement}&\textbf{Field}&\textbf{\(\dim X\)}&\textbf{Class}&
\textbf{Extra hypotheses}\\ \hline
\Cref{thm:bfj-envelope-regularity}
 &cdv, res.\ char.\ \(0\)
 &any
 &ample
 &\(X\) smooth projective; nothing further\\
\Cref{thm:ch23-gjkm-curves}
 &any complete, nontrivial
 &\(1\)
 &ample \(L\)
 &smooth projective geometrically connected curve; metric envelope\\
\Cref{thm:ch23-gjkm-positive-characteristic}
 &cdv, \(\operatorname{char}K=p\)
 &any
 &ample
 &geometric origin from a \(d\)-dimensional family over perfect \(k\);
   resolution over \(k\) in dimension \(d+n\)\\
\Cref{cor:ch23-gjkm-surfaces}
 &cdv, \(\operatorname{char}K=p\)
 &\(2\)
 &ample
 &geometric origin from a \(1\)-dimensional family; unconditional\\
\Cref{ex:ch23-outside-scope}
 &cdv, mixed char.
 &\(\geq3\)
 &ample
 &not covered by any of the above\\
\end{tabular}
\end{center}

Here ``cdv'' abbreviates ``complete discretely valued''.  None of these
statements licenses unrestricted continuity over every complete field, every
variety, and every numerical class.

\section{Chapter dictionary}
\label{sec:ch23-dictionary}

\subsection{From the objects of this chapter to their use}
\label{sec:ch23-dict-forward}

\begin{center}
\small
\begin{tabular}{p{.26\textwidth}|p{.36\textwidth}|p{.30\textwidth}}
\textbf{Object}&\textbf{Established property}&\textbf{First use}\\ \hline
obstacle envelope \(P_\theta(f)\)
 &sup-norm contraction, translation and reference equivariance
  (\Cref{prop:ch23-formal-envelope-calculus})
 &\Cref{lem:ch23-reduction-to-base-ideals},
  \Cref{chap:bfj-differentiability}\\
model envelope \(P^{\mathcal D}_\theta(f)\)
 &equals \(P_\theta(f)\) at quasimonomial points
  (\Cref{prop:ch23-model-vs-psh-envelope}), and everywhere after
  \Cref{cor:ch23-envelope-by-model-functions}
 &\Cref{sec:ch23-later}, \Cref{chap:bfj-continuity}\\
regularized supremum \((\sup_\alpha\varphi_\alpha)^*\)
 &\(\theta\)-psh; equals the raw sup at quasimonomial points
  (\Cref{thm:ch23-upper-envelope})
 &\Cref{thm:ch21-decreasing-stability},
  \Cref{chap:bfj-capacity}\\
continuous obstacle
 &raw envelope is already psh and maximal
  (\Cref{cor:ch23-continuous-obstacle-psh})
 &\Cref{thm:bfj-decreasing-regularization}\\
bounded obstacle
 &raw envelope may fail to be usc; no maximal competitor
  (\Cref{ex:ch23-bounded-obstacle-failure})
 &warning only\\
retraction of an obstacle
 &\(f\circ p_{\mathcal X}\to f\) uniformly
  (\Cref{lem:ch23-retraction-uniform})
 &\Cref{prop:ch23-model-vs-psh-envelope}\\
base ideal \(\mathfrak a_m\)
 &\(\varphi_m\) is \(\theta\)-psh, superadditive, and
  \(\varphi_m\to P_\theta(0)\) at rate \(O(m^{-1})\)
  (\Cref{thm:ch23-base-ideal-approximation},
   \Cref{cor:ch23-envelope-rate})
 &\Cref{thm:bfj-envelope-regularity}\\
asymptotic multiplier ideal \(\mathfrak b_m\)
 &containment, subadditivity, fixed-twist generation
  (\Cref{thm:ch23-uniform-generation-input})
 &Step~4 of
  \Cref{thm:ch23-base-ideal-approximation}\\
fixed-skeleton selection
 &increasing sequence of finite maxima, uniform on one \(\Delta\)
  (\Cref{prop:ch23-fixed-skeleton-selection})
 &\Cref{thm:bfj-supported-regularity}\\
model orthogonality, \(n=1\)
 &componentwise complementary slackness
  (\Cref{thm:ch23-model-orthogonality-curve})
 &\Cref{chap:bfj-examples}\\
convex orthogonality
 &\(\MA_{\mathbf R}\) does not charge the noncontact set
  (\Cref{thm:ch23-convex-envelope-orthogonality})
 &\Cref{chap:bfj-examples},
  \Cref{chap:compare-ma}\\
\end{tabular}
\end{center}

\subsection{The three envelopes, side by side}
\label{sec:ch23-dict-three}

\begin{center}
\small
\begin{tabularx}{\textwidth}{@{}>{\raggedright\arraybackslash}p{.12\textwidth}|%
 >{\raggedright\arraybackslash}X|>{\raggedright\arraybackslash}X|%
 >{\raggedright\arraybackslash}X@{}}
&\(P_\theta(f)\)&\(P^{\mathcal D}_\theta(f)\)
&\((\sup_\alpha\varphi_\alpha)^*\)\\ \hline
input&any \(f:X^{\An}\to\Rbar\)&any \(f\)&a psh family\\
competitors&\(\PSH(X,\theta)\)&\(\theta\)-psh model functions&none\\
usc for free?&only if \(f\) is usc
 (\Cref{cor:ch23-usc-obstacle})&no&yes, by construction\\
psh for free?&only if \(f\) is usc&no&yes
 (\Cref{thm:ch23-upper-envelope})\\
agreement&\multicolumn{2}{c|}{equal at quasimonomial points always,
 everywhere for continuous \(f\)}
&equals the raw sup at quasimonomial points\\
\end{tabularx}
\end{center}

\section*{Notes and further reading}

\enlargethispage{6\baselineskip}

\paragraph{The original theory.}
The whole of \Cref{sec:ch23-obstacle-operator} through
\Cref{sec:ch23-continuity} follows
Boucksom--Favre--Jonsson, \emph{Singular semipositive metrics in
non-Archimedean geometry} (2016) \cite{BFJ16}.  Upper-envelope stability is
\cite[Theorem~7.11]{BFJ16}; the formal calculus of
\Cref{prop:ch23-formal-envelope-calculus} is \cite[Proposition~8.2]{BFJ16};
continuity of the envelope of a continuous obstacle is
\cite[Theorem~8.3]{BFJ16}; base-ideal approximation is
\cite[Theorem~8.5]{BFJ16}; and monotone regularization, proved in
\Cref{chap:bfj-compactness} on the basis of the present chapter, is
\cite[Theorem~8.7]{BFJ16}.  The comparison with Zhang's continuous
semipositive metrics, \Cref{thm:ch21-continuous-zhang}, is
\cite[Corollary~8.8]{BFJ16}, and the identification of the two envelope
definitions on quasimonomial points, our
\Cref{prop:ch23-model-vs-psh-envelope}, is the content of
\cite[Lemma~8.4 and Lemma~8.9]{BFJ16}.  The multiplier-ideal package of
\Cref{thm:ch23-uniform-generation-input} is
\cite[Appendix~B, in particular Theorem~B.8 and Remark~B.9]{BFJ16}; its
absolute counterparts, together with the subadditivity theorem of
Demailly--Ein--Lazarsfeld \cite{DEL00} and the asymptotic formalism of
\cite{PAG2} and \cite{JM12}, are collected in \Cref{app:ideals}.

\paragraph{The presentational choices made here.}
Four things are organized differently from the source, and the reasons are
worth stating.  The two definitions of the envelope are separated at the
outset (\Cref{def:ch23-envelope} and \Cref{def:ch23-model-envelope}) and
their agreement is proved, because the later literature works with the
second and because the agreement is used in
\Cref{sec:ch23-later}.  The concavity estimate inside the proof of
\Cref{prop:ch23-formal-envelope-calculus}\,(7) is isolated as
\Cref{lem:ch23-concave-lipschitz}, since the version usually quoted omits
the passage to a concentric ball, which is where the constant comes from.
The uniform convergence of retractions,
\Cref{lem:ch23-retraction-uniform}, is stated once and used four times.  And
the model-level orthogonality of
\Cref{sec:ch23-model-orthogonality} is proved here, where the obstacle
problem first appears, rather than in
\Cref{chap:bfj-differentiability}, where the general statement lives.

\paragraph{The fixed-skeleton selection.}
The selection principle in \Cref{prop:ch23-fixed-skeleton-selection} is the
elementary dense-set, common-modulus, finite-maxima construction, followed
by Dini convergence.  The exact current statement-and-proof block was
verified in Rethlas run
\nolinkurl{20260824T044830Z_25024c43d265}, item
\nolinkurl{pi_4cc2365c9c77af25df553e8e}, with no critical issue or proof gap.
This is a book-local consequence assembled from the published skeletal
inputs, not a claim of publication priority or an attribution of the
proposition itself to the literature.  The
principle is deliberately restricted to one compact dual complex.  A global
sequential version---a
single sequence extracted from \(\mathcal F\) that works simultaneously on
all skeleta---would require a countably cofinal system of models and is
false in general; see \Cref{rem:ch23-selection-scope},
\Cref{rem:cofinality-scope}, and \Cref{ex:no-cofinal-sequence}.  The
statement is used, in exactly the fixed-skeleton form, in the
supported-measure regularity argument of
\Cref{thm:bfj-supported-regularity}.  Its relation to Choquet's lemma is
recorded in \Cref{rem:ch23-selection-vs-choquet}: the extra strength comes
from equicontinuity, not from any additional set theory.

\paragraph{Later continuity theorems, dated and sourced individually.}
Gubler--Jell--K\"unnemann--Martin, \emph{Continuity of plurisubharmonic
envelopes in non-Archimedean geometry and test ideals}, Annales de
l'Institut Fourier \textbf{69} (2019), 2331--2376, with an appendix by
Burgos Gil and Sombra \cite{GJKM19}, contains three results used above.
For ample line bundles on curves over an arbitrary complete nontrivially
valued field, continuity of the semipositive metric envelope is
\cite[Theorem~1.1]{GJKM19}, recorded in
\Cref{thm:ch23-gjkm-curves}.  No unrestricted degree-zero model-existence
assertion is used.  Continuity in equal
characteristic \(p\) is \cite[Theorem~8.2]{GJKM19}, our
\Cref{thm:ch23-gjkm-positive-characteristic}, with the unconditional
surface case \cite[Corollary~8.3]{GJKM19}.  The uniform generation theorem
for asymptotic test ideals is \cite[Theorem~1.3]{GJKM19}, proved over a base
as \cite[Theorem~4.6]{GJKM19}; it goes back to Mustaţă in the projective
case, and the relevant test-ideal formalism is that of
\cite{MustataTest}, \cite{HY03} and \cite{ST12}, reviewed in
\Cref{sec:appC-test-ideals}.  The Burgos Gil--Sombra appendix to
\cite{GJKM19} contains the toric counterexample recalled in
\Cref{rem:ch23-curve-specialness}: two models of a toric surface with the
same skeleton and different retractions, showing that the retraction of a
semipositive model metric need not be semipositive when \(n\geq2\).

The removal of algebraizability from the orthogonality property is
Burgos Gil--Gubler--Jell--K\"unnemann--Martin, \emph{Differentiability of
non-Archimedean volumes and non-Archimedean Monge--Amp\`ere equations},
Algebraic Geometry \textbf{7} (2020), 113--152 \cite{BGGJKM20}, with an
appendix by Lazarsfeld; the relevant statement is
\cite[Theorem~6.3.3]{BGGJKM20}, recorded here as
\Cref{thm:bggjkm-remove-algebraizability} and, with its exact provenance,
as \Cref{thm:orthogonality-provenance}.  Non-Archimedean volumes of nef
metrized line bundles are developed in \cite{BGM21}.  The
mixed-characteristic Monge--Amp\`ere theory currently available is
\cite{FGK25}; see \Cref{thm:fgk-mixed-characteristic-cy}.  Model metrics and
semipositivity over more general complete fields, which is the setting in
which \Cref{thm:ch23-gjkm-curves} is stated, are treated in
\cite{GM19}.  Alterations are \cite{deJong1996}; their exact role, and their
inability to substitute for resolution in the present argument, are
discussed in \Cref{sec:ch23-later-alterations}.

\paragraph{Orthogonality.}
Orthogonality in the sense of
\eqref{eq:ch23-orthogonality-preview} is not proved in this chapter for
\(n\geq2\).  It appeared first, under an algebraizability hypothesis, as
\cite[Theorem~A.4]{BFJ15}, and its later unconditional form is
\cite{BGGJKM20}; \Cref{thm:orthogonality-provenance} keeps the two apart.
What \emph{is} proved here is the model-level statement
\Cref{thm:ch23-model-orthogonality-curve} on a curve and the convex
statement \Cref{thm:ch23-convex-envelope-orthogonality}; both are
elementary, both are sharp in the sense that
\Cref{rem:ch23-model-slackness-fails} identifies exactly which of their
ingredients fails in higher dimension, and neither is used to prove
anything about continuity.  In particular this chapter proves continuity
without assuming orthogonality and never invokes a componentwise
complementary-slackness argument outside dimension one.  The convex
argument, through the nondifferentiability set of the Legendre conjugate, is
the same one that underlies \Cref{thm:appD-alexandrov-measure}; the
comparison of real and non-Archimedean Monge--Amp\`ere operators in the
toric case is \cite{Vil20}, and its global invariant tropical form is
\cite{BGJKK25}.

\paragraph{What is new here, and what is not.}
No theorem of this chapter is new.
\Cref{thm:ch23-model-orthogonality-curve} and
\Cref{thm:ch23-convex-envelope-orthogonality} are folklore statements of
linear programming and convex analysis respectively, written out because
\Cref{rem:model-slackness-heuristic} refers to them and because the
worked computations of \Cref{sec:ch23-worked} are checked against them.
\Cref{ex:ch23-outside-scope} is a scope statement, not a counterexample: it
records that continuity of envelopes for threefolds over a
mixed-characteristic complete discretely valued field is not covered by any
theorem proved or cited in this book, and it should not be read as an
assertion that continuity fails there.
